\documentclass[11pt]{article}
\usepackage[margin=1in]{geometry}
\usepackage[T1]{fontenc}
\usepackage[utf8]{inputenc}
\usepackage[french]{babel}
\usepackage{amsmath,amssymb,amsthm}
\usepackage[colorlinks=true,linkcolor=blue,urlcolor=blue]{hyperref}
\newtheorem{problem}{Problème}
\newenvironment{refs}{\par\small\noindent\emph{Références :}\begin{itemize}\setlength\itemsep{0pt}}{\end{itemize}\normalsize}
\title{Hors de la Caverne \\ \Large Équations différentielles ordinaires et systèmes dynamiques}
\author{}
\date{}
\begin{document}
\maketitle
\noindent\emph{Des problèmes, pas de réponses.} Chaque problème ci-dessous est posé
sans solution. Les références indiquent où trouver les connaissances nécessaires.
\medskip


\section*{Lycée}

\begin{problem}[{La croissance exponentielle à partir des principes premiers --- Lycée}]
\textbf{ODE-01.} \textbf{Problème.} (a) Démontrer que les \textit{seules} fonctions dérivables
vérifiant $ y' = k y $ sur la droite réelle sont $ y(x) = C e^{kx} $ avec $ C $ constante.
(On pourra utiliser les règles élémentaires de dérivation et le fait qu'une fonction de dérivée
nulle sur un intervalle y est constante.)
(b) En déduire la loi du temps de doublement : une solution positive avec $ k >0 $ double en
un même temps $ (\ln 2)/k $, quel que soit l'instant où l'on déclenche le chronomètre.
(c) Datation au radiocarbone : le carbone 14 se désintègre avec une demi-vie de 5730 ans. Un
parchemin conserve 70 \% de son carbone 14 initial. Déterminer son âge, en le justifiant
entièrement à partir de (a).
\end{problem}

\noindent C'est la première équation différentielle jamais comprise --- elle régit les
intérêts composés (Jacques Bernoulli, années 1690), le refroidissement, la désintégration
radioactive (Rutherford et Soddy, 1902) et la croissance des populations (Malthus, 1798). La
démonstration d'unicité en (a) est la partie importante : c'est le cas le plus simple du fait
qu'une équation différentielle assortie d'une condition initiale détermine complètement l'avenir,
le thème de toute cette page. Willard Libby a reçu le prix Nobel de chimie 1960 pour avoir fait
de (c) un instrument d'archéologie.

\begin{refs}
\item Contexte : Croissance exponentielle --- Wikipédia (\url{https://fr.wikipedia.org/wiki/Croissance_exponentielle})
\item Lecture : W. E. Boyce \& R. C. DiPrima, \textit{Elementary Differential Equations}, chap. 1--2
\item Cours : MIT OCW 18.03SC Differential Equations (\url{https://ocw.mit.edu/courses/18-03sc-differential-equations-fall-2011/})
\end{refs}

\bigskip

\begin{problem}[{La loi de refroidissement de Newton, à partir de données --- Lycée}]
\textbf{ODE-02.} \textbf{Problème.} La loi de refroidissement de Newton affirme que la température $ T $ d'un
corps placé dans un milieu à température constante $ A $ obéit à $ T' = -k (T - A) $ pour une
certaine constante $ k >0 $.
\begin{enumerate}
\item[(a)] Résoudre l'équation (utiliser ODE-01 après la substitution $ u = T - A $).
\item[(b)] Une tasse de café à 90$^\circ$C se trouve dans une pièce à 20$^\circ$C. Elle mesure 70$^\circ$C au
bout de 5 minutes et environ 55,7$^\circ$C au bout de 10 minutes. Montrer que ces données sont
compatibles avec le modèle, déterminer $ k $, puis prédire la température au bout de 20 minutes
ainsi que l'instant où le café atteint 40$^\circ$C.
\item[(c)] Démontrer que, selon le modèle, le café n'atteint jamais réellement la température de la
pièce en temps fini --- et expliquer pourquoi le modèle est néanmoins utile.
\end{enumerate}

\end{problem}

\noindent Newton a publié la loi de refroidissement anonymement en 1701, dans un article
sur une échelle de chaleur allant jusqu'à la fusion du fer ; confronter un modèle d'équation
différentielle à des données, comme en (b), est exactement ce qu'il faisait. Cette loi est le
premier exemple canonique d'ajustement d'une équation différentielle à des mesures --- les
médecins légistes s'en servent encore pour estimer l'heure du décès --- et la partie (c) apprend
la bonne façon de lire une approche exponentielle de l'équilibre.

\begin{refs}
\item Contexte : Loi de refroidissement de Newton --- Wikipédia (\url{https://fr.wikipedia.org/wiki/Loi_de_refroidissement_de_Newton})
\item Lecture : W. E. Boyce \& R. C. DiPrima, \textit{Elementary Differential Equations}, chap. 1
\item Cours : MIT OCW 18.03SC Differential Equations (\url{https://ocw.mit.edu/courses/18-03sc-differential-equations-fall-2011/})
\end{refs}

\bigskip

\begin{problem}[{La courbe logistique --- Lycée}]
\textbf{ODE-03.} \textbf{Problème.} Une population limitée par les ressources est modélisée par l'équation
logistique $ P' = r P \left( 1 - \frac{P}{K} \right) $ avec des constantes
$ r, K >0 $.
\begin{enumerate}
\item[(a)] Sans résoudre, montrer que toute solution avec $ 0 <P(0) <K $ est croissante, et que
toute solution avec $ P(0) >K $ est décroissante.
\item[(b)] Résoudre l'équation par séparation des variables et décomposition en éléments simples, puis
vérifier que $ P(t) \to K $ quand $ t \to \infty $ pour toute valeur initiale
positive.
\item[(c)] Démontrer qu'une solution partant en dessous de $ K/2 $ a sa croissance la plus forte
exactement lorsque $ P = K/2 $, et que la courbe solution est symétrique par rapport à ce
point. Esquisser la courbe en S.
\end{enumerate}

\end{problem}

\noindent Le mathématicien belge Pierre-François Verhulst a proposé l'équation logistique
en 1838 pour corriger Malthus : la croissance ne peut pas rester exponentielle. La même courbe en
S modélise aujourd'hui les épidémies, l'adoption des technologies et l'autocatalyse chimique ;
sa cousine à temps discret reviendra en ODE-16, chargée de chaos. Verhulst s'en est servi pour
prédire un plafond à la population belge --- un rappel que les paramètres du modèle, contrairement
à ses mathématiques, sont des conjectures.

\begin{refs}
\item Contexte : Fonction logistique (Verhulst) --- Wikipédia (\url{https://fr.wikipedia.org/wiki/Fonction_logistique_(Verhulst)})
\item Lecture : S. H. Strogatz, \textit{Nonlinear Dynamics and Chaos}, chap. 2
\item Cours : MIT OCW 18.03SC Differential Equations (\url{https://ocw.mit.edu/courses/18-03sc-differential-equations-fall-2011/})
\end{refs}

\bigskip

\begin{problem}[{La vitesse limite --- Lycée, short}]
\textbf{ODE-22.} \textbf{Problème.} Un corps de masse $ m $ tombe sans vitesse initiale sous
l'effet de la pesanteur, contre une résistance de l'air proportionnelle à la vitesse, de sorte
que sa vitesse vérifie
\[ m v' = mg - k v, \qquad v(0) = 0, \]
avec des constantes $ m, g, k >0 $. Résoudre en $ v(t) $ et démontrer que $ v $ croît
strictement vers la vitesse limite $ mg/k $ sans jamais l'atteindre.
\end{problem}

\noindent Le frottement linéaire est le modèle honnête d'une chute lente --- une
gouttelette de brume, une bille d'acier dans l'huile --- et Stokes en a établi la loi en 1851. La
moitié qualitative de ce problème est la plus précieuse : la limite $ mg/k $ se lit sur
l'équation avant toute intégration, en demandant où $ v' = 0 $, ce qui est le premier geste de
toute la théorie qualitative vers laquelle cette page tend.

\begin{refs}
\item Contexte : Vitesse terminale --- Wikipédia (\url{https://fr.wikipedia.org/wiki/Vitesse_terminale})
\item Lecture : W. E. Boyce \& R. C. DiPrima, \textit{Elementary Differential Equations}, chap. 1
\end{refs}

\bigskip

\begin{problem}[{Le réservoir qui se vide en temps fini --- Lycée, short}]
\textbf{ODE-23.} \textbf{Problème.} De l'eau s'écoule par un petit trou percé au fond d'un
réservoir cylindrique ; la loi de Torricelli impose à la hauteur $ h(t) $ d'obéir à
\[ h' = -c\sqrt{h}, \qquad h(0) = h_0 >0, \]
avec une constante $ c >0 $ déterminée par le trou. Déterminer, avec démonstration, l'instant
exact où le réservoir est vide, et démontrer que la hauteur est véritablement nulle à partir de
cet instant, et non pas seulement tendant vers zéro.
\end{problem}

\noindent Evangelista Torricelli --- élève de Galilée et inventeur du baromètre --- a publié
la loi en 1644 : l'eau quitte le trou à la vitesse qu'elle atteindrait en tombant librement depuis
la surface. Le temps de vidange fini est tout l'intérêt. Les solutions de $ y' = y $ mettent
l'éternité à atteindre quoi que ce soit, mais $ \sqrt{h} $ n'est pas lipschitzienne en $ 0 $,
et c'est ce défaut qui permet à une solution d'arriver à l'origine en temps fini --- le même défaut
qui détruit l'unicité en ODE-08.

\begin{refs}
\item Contexte : Formule de Torricelli --- Wikipédia (\url{https://fr.wikipedia.org/wiki/Formule_de_Torricelli})
\item Cours : MIT OCW 18.03SC Differential Equations (\url{https://ocw.mit.edu/courses/18-03sc-differential-equations-fall-2011/})
\end{refs}

\bigskip

\begin{problem}[{Champs de tangentes et méthode d'Euler --- Lycée}]
\textbf{ODE-24.} \textbf{Problème.} On considère le problème de Cauchy
\[ y' = y - t, \qquad y(0) = \tfrac{1}{2}. \]
\begin{enumerate}
\item[(a)] Tracer le champ de tangentes sur le rectangle $ 0 \le t \le 2 $, $ -1 \le y \le 3 $.
Vérifier que $ y = t + 1 $ est solution, et utiliser le signe de $ y - t $ au-dessus et
en dessous de cette droite pour prédire l'allure de toutes les autres solutions.
\item[(b)] Vérifier que la solution du problème de Cauchy est
$ y(t) = t + 1 - \tfrac{1}{2}e^{t} $, et confirmer qu'elle correspond à votre prédiction.
\item[(c)] Appliquer la méthode d'Euler $ y_{n+1} = y_n + h\,(y_n - t_n) $ avec les pas
$ h = 0{,}5,\ 0{,}25,\ 0{,}125 $ pour estimer $ y(1) $. Dresser le tableau des trois erreurs
et observer que chaque division par deux de $ h $ divise à peu près l'erreur par deux.
\item[(d)] Expliquer cette observation. Montrer par un développement de Taylor qu'un pas d'Euler à
partir d'une valeur exacte commet une erreur au plus $ \tfrac{1}{2}h^2 \max|y''| $, et
argumenter qu'en cumulant $ 1/h $ pas de ce type on obtient une erreur totale d'ordre $ h $ ---
la méthode d'Euler est donc du \textit{premier ordre}.
\end{enumerate}

\end{problem}

\noindent Leonhard Euler a publié cette méthode dans ses \textit{Institutionum calculi
integralis} (1768) comme moyen de calculer avec des équations que personne ne savait résoudre.
Elle est l'ancêtre de tous les solveurs numériques utilisés aujourd'hui, et aussi un instrument
théorique : Cauchy dans les années 1820, et Peano en 1886, ont démontré que les solutions
\textit{existent} en montrant que les approximations polygonales d'Euler convergent quand
$ h \to 0 $ (comparer à la voie différente d'ODE-07). Le champ de tangentes de (a) est l'autre
moitié de la leçon --- une équation qu'on ne sait pas résoudre peut tout de même se lire.

\begin{refs}
\item Contexte : Méthode d'Euler --- Wikipédia (\url{https://fr.wikipedia.org/wiki/M%C3%A9thode_d%27Euler})
\item Contexte : Champ de tangentes --- Wikipédia (en anglais) (\url{https://en.wikipedia.org/wiki/Slope_field})
\item Lecture : W. E. Boyce \& R. C. DiPrima, \textit{Elementary Differential Equations}, chap. 1--2
\item Cours : MIT OCW 18.03SC Differential Equations (\url{https://ocw.mit.edu/courses/18-03sc-differential-equations-fall-2011/})
\end{refs}

\bigskip

\begin{problem}[{La chaîne suspendue --- Lycée}]
\textbf{ODE-34.} \textbf{Problème.} Une chaîne souple homogène de poids $ w $ par unité de longueur pend au
repos entre deux supports. Soit $ y(x) $ sa forme ; on place le point le plus bas en
$ x = 0 $, et l'on note $ H $ la traction horizontale en ce point.
\begin{enumerate}
\item[(a)] Écrire l'équilibre des forces sur le morceau de chaîne compris entre le point le plus bas
et un point courant. Montrer que la composante horizontale de la tension vaut $ H $ en tout
point, que la composante verticale vaut $ w\,s $ où $ s $ est l'abscisse curviligne mesurée
depuis le point le plus bas, et donc que $ y'(x) = w\,s(x)/H $. Dériver, en utilisant
$ \frac{ds}{dx} = \sqrt{1 + (y')^2} $, pour obtenir
\[ y'' = \frac{w}{H}\,\sqrt{1 + (y')^2}\,. \]
\item[(b)] Résoudre l'équation : poser $ p = y' $, séparer les variables, et montrer que la chaîne
telle que $ y'(0) = 0 $ et $ y(0) = H/w $ est
\[ y(x) = \frac{H}{w}\,\cosh\!\left( \frac{w x}{H} \right). \]
Confirmer par dérivation directe que cette fonction satisfait l'équation.
\item[(c)] Donner tort à Galilée : montrer qu'aucune parabole $ y = ax^2 + bx + c $ avec $ a \ne 0 $
ne peut satisfaire l'équation, en comparant les deux membres. Puis développer
$ \cosh u = 1 + \frac{u^2}{2} + \frac{u^4}{24} + \cdots $ et expliquer pourquoi une chaîne peu
tendue paraît néanmoins parabolique.
\item[(d)] Changer le chargement : suspendre un tablier au câble, de sorte que la charge soit $ w_0 $
par unité de longueur \textit{horizontale}, le poids propre du câble étant négligeable. Refaire
l'équilibre des forces, établir la nouvelle équation, et démontrer que ce câble est exactement
une parabole. Laquelle de ces deux courbes suit le câble porteur d'un pont suspendu, et laquelle
suit un collier ?
\end{enumerate}

\end{problem}

\noindent Galilée affirmait dans les \textit{Discours concernant deux sciences nouvelles}
(1638) qu'une chaîne suspendue est une parabole ; Joachim Jungius a trouvé l'affirmation fausse,
et en 1690 Jacques Bernoulli a publiquement mis les mathématiciens d'Europe au défi de produire
la vraie courbe. Huygens, Leibniz et Jean Bernoulli ont chacun publié une solution dans les
\textit{Acta Eruditorum} en 1691 --- Huygens a forgé le nom de \textit{catenaria}, du latin
\textit{catena}, chaîne --- et la réponse, qui exige le cosinus hyperbolique, fut une victoire
précoce du calcul nouveau sur un problème que les Anciens ne pouvaient que deviner. Robert Hooke
avait vu la conséquence pour l'ingénierie dès 1675 et l'avait publiée sous forme d'anagramme
latine : comme pend un câble souple, ainsi, renversés, se tiennent les claveaux d'une arche. La
partie (d) explique pourquoi les deux formes sont sans cesse confondues --- c'est la charge, non la
courbe, qui diffère.

\begin{refs}
\item Contexte : Chaînette --- Wikipédia (\url{https://fr.wikipedia.org/wiki/Cha%C3%AEnette})
\item Contexte : Fonction hyperbolique --- Wikipédia (\url{https://fr.wikipedia.org/wiki/Fonction_hyperbolique})
\item Lecture : G. F. Simmons, \textit{Differential Equations with Applications and Historical Notes}, chap. 1
\item Cours : MIT OCW 18.03SC Differential Equations (\url{https://ocw.mit.edu/courses/18-03sc-differential-equations-fall-2011/})
\end{refs}

\bigskip

\begin{problem}[{La vitesse de libération --- Lycée, short}]
\textbf{ODE-35.} \textbf{Problème.} Un projectile est tiré verticalement depuis la surface
d'une planète sans atmosphère, de masse $ M $ et de rayon $ R $, à la vitesse
$ v_0 > 0 $, puis se déplace sous la seule action de la gravitation :
\[ \frac{d^2 r}{dt^2} = -\frac{GM}{r^2}, \qquad r(0) = R, \qquad r'(0) = v_0 . \]
En posant $ v = dr/dt $ et en utilisant $ \frac{dv}{dt} = v\,\frac{dv}{dr} $, déterminer avec
démonstration la plus petite valeur de $ v_0 $ pour laquelle $ r(t) $ croît sans borne, et
démontrer que pour toute valeur plus petite de $ v_0 $ le projectile atteint une hauteur
maximale puis retombe.
\end{problem}

\noindent Newton a dessiné l'image dans son \textit{Traité du système du monde} : un canon
sur une très haute montagne, tirant de plus en plus fort, jusqu'à ce que le boulet ne retombe
plus. La substitution demandée ici est toute l'astuce --- elle transforme une équation du second
ordre en une équation du premier ordre en la variable $ r $, et ce qu'elle produit est la
conservation de l'énergie, la même intégrale première qui donne la période du pendule en ODE-12.
Le nombre qu'elle fournit vaut environ 11,2 km/s pour la Terre et 2,4 km/s pour la Lune, ce qui
explique en grande partie pourquoi la Lune n'a pas d'atmosphère à perdre.

\begin{refs}
\item Contexte : Vitesse de libération --- Wikipédia (\url{https://fr.wikipedia.org/wiki/Vitesse_de_lib%C3%A9ration})
\item Contexte : Canon de Newton --- Wikipédia (\url{https://fr.wikipedia.org/wiki/Canon_de_Newton})
\item Lecture : W. E. Boyce \& R. C. DiPrima, \textit{Elementary Differential Equations}, chap. 1--2
\end{refs}

\bigskip


\section*{Licence}

\begin{problem}[{Équations du premier ordre : la cuve de mélange --- Licence}]
\textbf{ODE-04.} \textbf{Problème.} Une cuve contient 100 litres d'eau pure. De la saumure contenant 0,5 kg de
sel par litre y entre à 3 litres par minute ; la cuve est parfaitement mélangée et se vide à
3 litres par minute. Soit $ S(t) $ la masse de sel dans la cuve, en kilogrammes.
\begin{enumerate}
\item[(a)] Établir, en justifiant chaque terme à partir des principes premiers (débit entrant moins
débit sortant), l'équation $ S' = 1{,}5 - \frac{3S}{100} $.
\item[(b)] La résoudre par facteur intégrant, et démontrer que $ S(t) \to 50 $ de façon
monotone.
\item[(c)] Supposer maintenant que le débit sortant vaut 2 litres par minute, de sorte que le volume
augmente. Établir la nouvelle équation, à coefficient dépendant du temps, et la résoudre sur
l'intervalle précédant le débordement de la cuve.
\item[(d)] Expliquer pourquoi toute équation linéaire du premier ordre $ y' + p(t) y = q(t) $ se
plie à la même méthode du facteur intégrant, et démontrer que cette méthode donne
\textit{toutes} les solutions.
\end{enumerate}

\end{problem}

\noindent Les problèmes de mélange sont le terrain d'entraînement classique des équations
linéaires du premier ordre --- la même loi de bilan régit les polluants dans les lacs, les
médicaments dans le sang et le CO$_2$ dans une pièce. Le facteur intégrant remonte à Leibniz et au
traitement systématique des équations du premier ordre par Euler dans les années 1730--1740.
L'intérêt de (a) est la discipline de modélisation : chaque terme de l'équation doit être un
débit physique dûment justifié.

\begin{refs}
\item Contexte : Équation différentielle linéaire --- Wikipédia (\url{https://fr.wikipedia.org/wiki/%C3%89quation_diff%C3%A9rentielle_lin%C3%A9aire})
\item Lecture : W. E. Boyce \& R. C. DiPrima, \textit{Elementary Differential Equations}, chap. 2
\item Cours : MIT OCW 18.03 Differential Equations (\url{https://ocw.mit.edu/courses/18-03-differential-equations-spring-2010/})
\end{refs}

\bigskip

\begin{problem}[{L'oscillateur harmonique, entièrement résolu --- Licence}]
\textbf{ODE-05.} \textbf{Problème.} On considère $ m y'' + c y' + k y = 0 $ avec
$ m, k >0 $ et un amortissement $ c \ge 0 $.
\begin{enumerate}
\item[(a)] Montrer que la recherche de solutions $ e^{rt} $ conduit à l'équation caractéristique
$ m r^2 + c r + k = 0 $, et écrire la solution générale dans les trois cas (racines réelles
distinctes, racine double, racines complexes), en \textit{établissant} la solution $ t e^{rt} $
dans le cas de la racine double plutôt qu'en la devinant.
\item[(b)] Démontrer que ce sont bien là \textit{toutes} les solutions. (Utiliser le théorème d'existence
et d'unicité pour les équations linéaires du second ordre, ou démontrer directement qu'une
solution s'annulant en un point ainsi que sa dérivée est identiquement nulle.)
(c) Démontrer : toute solution non nulle tend vers 0 quand $ t \to \infty $ si et seulement si
$ c >0 $. Interpréter physiquement les trois régimes d'amortissement, et expliquer pourquoi
la suspension d'une voiture est conçue près de l'amortissement critique.
\end{enumerate}

\end{problem}

\noindent L'oscillateur est la drosophile des équations différentielles : ressorts,
pendules de petite amplitude, circuits RLC et vibrations moléculaires sont tous cette unique
équation. Euler a relié l'exponentielle aux sinus et cosinus dans les années 1740, et la méthode
de l'équation caractéristique est essentiellement la sienne. La partie (b) est l'endroit où le
sujet devient honnête --- une formule n'est pas un théorème tant qu'on ne sait pas qu'aucune
solution ne manque.

\begin{refs}
\item Contexte : Oscillateur harmonique --- Wikipédia (\url{https://fr.wikipedia.org/wiki/Oscillateur_harmonique})
\item Lecture : W. E. Boyce \& R. C. DiPrima, \textit{Elementary Differential Equations}, chap. 3
\item Cours : MIT OCW 18.03 Differential Equations (\url{https://ocw.mit.edu/courses/18-03-differential-equations-spring-2010/})
\end{refs}

\bigskip

\begin{problem}[{La résonance, rigoureusement --- Licence}]
\textbf{ODE-06.} \textbf{Problème.} On excite un oscillateur non amorti à la pulsation $ \omega $ :
$ x'' + \omega_0^2 x = F \cos(\omega t) $, avec $ F \ne 0 $.
\begin{enumerate}
\item[(a)] Pour $ \omega \ne \omega_0 $, trouver la solution particulière
$ \frac{F}{\omega_0^2 - \omega^2} \cos(\omega t) $ et démontrer que toute solution est bornée.
En partant du repos, montrer que la solution est une superposition de deux cosinus et --- lorsque
$ \omega $ est proche de $ \omega_0 $ --- la réécrire comme un produit faisant apparaître des
\textit{battements}.
\item[(b)] Pour $ \omega = \omega_0 $, démontrer que $ \frac{F}{2\omega_0}\, t \sin(\omega_0 t) $
est solution de l'équation et que \textit{toute} solution est non bornée : la résonance n'est pas
une possibilité mais un théorème.
\item[(c)] Ajouter l'amortissement : $ x'' + c x' + \omega_0^2 x = F \cos(\omega t) $ avec
$ c >0 $. Montrer que toutes les solutions convergent vers une unique réponse périodique
d'amplitude $ F / \sqrt{ (\omega_0^2 - \omega^2)^2 + c^2 \omega^2 } $, et déterminer pour
quelles valeurs de $ c $ cette amplitude admet un véritable maximum en $ \omega $ (le pic de
résonance) et où il se situe.
\end{enumerate}

\end{problem}

\noindent La résonance explique pourquoi les soldats rompent le pas sur les ponts et
pourquoi un verre à vin peut être brisé par la voix. L'effondrement du pont de Tacoma Narrows en
1940 en est l'illustration classique --- bien que, comme Billah et Scanlan l'ont montré dans un
article célèbre de 1991, le vrai coupable fût le flottement aéroélastique, une oscillation
auto-entretenue, et non une simple résonance forcée ; apprendre ce sujet, c'est aussi apprendre à
énoncer honnêtement les exemples. Les mathématiques de (b), le terme séculaire croissant comme
$ t $, sont exactement ce que les théoriciens des perturbations en mécanique céleste ont
combattu pendant deux siècles.

\begin{refs}
\item Contexte : Résonance --- Wikipédia (\url{https://fr.wikipedia.org/wiki/R%C3%A9sonance})
\item Lecture : W. E. Boyce \& R. C. DiPrima, \textit{Elementary Differential Equations}, chap. 3
\item Article : K. Y. Billah \& R. H. Scanlan, \og{} Resonance, Tacoma Narrows bridge failure, and undergraduate physics textbooks \fg{}, \textit{American Journal of Physics} 59 (1991), 118--124
\end{refs}

\bigskip

\begin{problem}[{Cauchy--Lipschitz (Picard--Lindelöf) : existence et unicité --- Licence}]
\textbf{ODE-07.} \textbf{Problème.} On considère le problème de Cauchy $ y' = f(t, y) $,
$ y(t_0) = y_0 $, où $ f $ est continue et lipschitzienne en $ y $ :
$ |f(t, y_1) - f(t, y_2)| \le L |y_1 - y_2| $.
\begin{enumerate}
\item[(a)] Montrer que le problème équivaut à l'équation intégrale
$ y(t) = y_0 + \int_{t_0}^{t} f(s, y(s))\, ds $.
\item[(b)] Démontrer que, sur un intervalle assez court, l'application qui envoie une fonction continue
$ \phi $ sur $ t \mapsto y_0 + \int_{t_0}^{t} f(s, \phi(s))\, ds $ est contractante pour la
norme de la convergence uniforme, et conclure du théorème du point fixe de Banach que le problème
de Cauchy y admet exactement une solution.
\item[(c)] Faire tourner l'itération à la main : pour $ y' = y $, $ y(0) = 1 $, en partant de la
fonction constante 1, calculer les itérées de Picard et identifier la limite.
\end{enumerate}

\end{problem}

\noindent Qu'une équation différentielle assortie d'un état initial détermine une
trajectoire unique est la forme mathématique du déterminisme, et c'est un théorème, non une
évidence --- démontré par Émile Picard (1890) et Ernst Lindelöf (1894) par l'itération que vous
menez ici. L'argument fut l'un des premiers triomphes du point de vue du point fixe, cristallisé
plus tard par Banach (1922), et la restriction à un intervalle court est bien réelle : les
solutions d'équations aussi innocentes que $ y' = y^2 $ explosent en temps fini.

\begin{refs}
\item Contexte : Théorème de Cauchy-Lipschitz --- Wikipédia (\url{https://fr.wikipedia.org/wiki/Th%C3%A9or%C3%A8me_de_Cauchy-Lipschitz})
\item Lecture : M. W. Hirsch, S. Smale \& R. L. Devaney, \textit{Differential Equations, Dynamical Systems, and an Introduction to Chaos}, chap. 17
\item Lecture : V. I. Arnold, \textit{Ordinary Differential Equations}
\end{refs}

\bigskip

\begin{problem}[{Là où l'unicité tombe en défaut --- Licence}]
\textbf{ODE-08.} \textbf{Problème.} On considère $ y' = y^{2/3} $ avec $ y(0) = 0 $.
\begin{enumerate}
\item[(a)] Exhiber deux solutions distinctes passant par l'origine.
\item[(b)] Bien pire : construire une famille infinie de solutions --- des fonctions qui restent nulles
jusqu'à un instant $ a \ge 0 $ choisi arbitrairement et ne décollent qu'ensuite --- et vérifier
soigneusement que chacune est dérivable partout et satisfait l'équation partout, y compris au
point de raccordement.
\item[(c)] Identifier exactement quelle hypothèse du théorème de Cauchy--Lipschitz (ODE-07) est mise en
défaut en $ y = 0 $, et démontrer que par toute valeur initiale $ y_0 \ne 0 $ la solution
\textit{est} localement unique.
\item[(d)] Énoncer le théorème d'existence de Peano et expliquer ce que la seule continuité de $ f $
vous achète et ce qu'elle ne vous achète pas.
\end{enumerate}

\end{problem}

\noindent Cette équation d'apparence innocente est le témoin classique du fait que le
déterminisme peut échouer lorsque le second membre est seulement continu : le seau percé qui
pourrait s'être vidé à n'importe quel moment du passé, vu à rebours. Peano a démontré l'existence
sans l'unicité en 1890 ; l'exemple montre que son théorème est optimal. Des modèles physiquement
sensés à second membre non lipschitzien existent bel et bien, et savoir exactement où la garantie
classique se rompt fait partie de la boîte à outils de l'analyste.

\begin{refs}
\item Contexte : Théorème de Cauchy-Peano-Arzelà --- Wikipédia (\url{https://fr.wikipedia.org/wiki/Th%C3%A9or%C3%A8me_de_Cauchy-Peano-Arzel%C3%A0})
\item Contexte : Théorème de Cauchy-Lipschitz --- Wikipédia (\url{https://fr.wikipedia.org/wiki/Th%C3%A9or%C3%A8me_de_Cauchy-Lipschitz})
\item Lecture : V. I. Arnold, \textit{Ordinary Differential Equations}
\end{refs}

\bigskip

\begin{problem}[{Le plan de phase : classification des systèmes linéaires --- Licence}]
\textbf{ODE-09.} \textbf{Problème.} On considère le système linéaire plan $ x' = A x $ pour une matrice réelle
$ 2 \times 2 $ $ A $ avec $ \det A \ne 0 $.
\begin{enumerate}
\item[(a)] Résoudre le système dans chaque cas de valeurs propres (réelles distinctes, complexes
conjuguées, double) et esquisser les portraits de phase : col, nœuds, foyers, centres.
\item[(b)] Démontrer la classification par trace et déterminant : l'origine est un col si et seulement
si $ \det A <0 $ ; elle est asymptotiquement stable si et seulement si
$ \operatorname{tr} A <0 $ et $ \det A >0 $ ; les foyers apparaissent exactement quand
$ (\operatorname{tr} A)^2 <4 \det A $. Dessiner le plan trace--déterminant avec toutes les
régions étiquetées.
\item[(c)] Expliquer pourquoi le cas du centre (valeurs propres imaginaires pures) est \textit{fragile} :
exhiber des matrices arbitrairement proches d'un centre dont les portraits sont des foyers
attractifs et des foyers répulsifs.
\end{enumerate}

\end{problem}

\noindent Le portrait de phase --- les trajectoires tracées dans l'espace des états, le
temps restant implicite --- est la grande invention picturale de Poincaré (1881--1886), et la
classification linéaire donnée ici est l'alphabet de la théorie qualitative : au voisinage d'un
équilibre typique, un système non linéaire ressemble à l'une de ces images (ODE-10). La fragilité
de (c) est le premier indice d'un thème profond, la stabilité structurelle : certains portraits
de phase survivent aux perturbations, d'autres non.

\begin{refs}
\item Contexte : Portrait de phase --- Wikipédia (\url{https://fr.wikipedia.org/wiki/Portrait_de_phase})
\item Lecture : S. H. Strogatz, \textit{Nonlinear Dynamics and Chaos}, chap. 5
\item Lecture : M. W. Hirsch, S. Smale \& R. L. Devaney, \textit{Differential Equations, Dynamical Systems, and an Introduction to Chaos}, chap. 2--4
\end{refs}

\bigskip

\begin{problem}[{La linéarisation et ses limites --- Licence}]
\textbf{ODE-10.} \textbf{Problème.} Soit $ x^* $ un point d'équilibre du système non linéaire
$ x' = f(x) $ dans le plan, et soit $ J $ la matrice jacobienne de $ f $ en $ x^* $.
\begin{enumerate}
\item[(a)] Énoncer précisément le théorème de Hartman--Grobman (équilibre hyperbolique : aucune valeur
propre de $ J $ sur l'axe imaginaire), et s'en servir pour classer tous les équilibres du
pendule amorti $ \theta'' + b \theta' + \sin\theta = 0 $ pour $ b >0 $ : montrer que les
positions basses sont des foyers ou des nœuds stables et que les positions renversées sont des
cols.
\item[(b)] Démontrer que la linéarisation peut mentir lorsque l'équilibre n'est pas hyperbolique : pour
le système $ x' = -y + a x (x^2 + y^2) $, $ y' = x + a y (x^2 + y^2) $, montrer que la
linéarisation à l'origine est un centre pour toute valeur de $ a $, mais qu'en calculant la
dérivée de $ r^2 = x^2 + y^2 $ le long des solutions on démontre que l'origine est
asymptotiquement stable pour $ a <0 $ et instable pour $ a >0 $.
\end{enumerate}

\end{problem}

\noindent La linéarisation est le premier geste du scientifique en exercice devant tout
équilibre, et Hartman--Grobman (démontré vers 1959--1960 par Philip Hartman et David Grobman
indépendamment) en est la licence. La partie (b) est l'exemple d'avertissement classique, chez
Strogatz et dans tout bon cours : sur le fil du rasoir des valeurs propres imaginaires, ce sont
les termes non linéaires qui décident, et fixer $ J $ du regard ne vous l'apprendra jamais. Le
problème centre-foyer de Poincaré --- décider lequel des deux --- est encore un sujet de recherche
pour les systèmes polynomiaux (voir ODE-19).

\begin{refs}
\item Contexte : Théorème de Hartman-Grobman --- Wikipédia (\url{https://fr.wikipedia.org/wiki/Th%C3%A9or%C3%A8me_de_Hartman-Grobman})
\item Lecture : S. H. Strogatz, \textit{Nonlinear Dynamics and Chaos}, chap. 6
\item Cours : MIT OCW 18.385J Nonlinear Dynamics and Chaos (\url{https://ocw.mit.edu/courses/18-385j-nonlinear-dynamics-and-chaos-fall-2014/})
\end{refs}

\bigskip

\begin{problem}[{Lotka--Volterra : les orbites sont fermées --- Licence}]
\textbf{ODE-11.} \textbf{Problème.} Le système proie--prédateur de Lotka et Volterra, pour des
proies $ x >0 $ et des prédateurs $ y >0 $ avec des paramètres positifs :
\[ x' = x (a - b y), \qquad y' = y (-c + d x) . \]
(a) Trouver le point d'équilibre dans le quadrant ouvert et montrer que la linéarisation y est un
centre --- donc ODE-10 avertit que la linéarisation seule ne décide de rien.
(b) Démontrer que $ H(x, y) = d x - c \ln x + b y - a \ln y $ est constante le long de toute
solution.
(c) Démontrer que chaque ligne de niveau de $ H $ dans le quadrant ouvert est une courbe fermée
bornée entourant le point d'équilibre, et en conclure que toute orbite non stationnaire est
périodique.
(d) Loi de Volterra : démontrer que la moyenne de $ x $ sur une période vaut la valeur
d'équilibre $ c/d $, indépendamment de l'orbite. (Intégrer $ (\ln y)' $ sur une
période.)
\end{problem}

\noindent Vito Volterra a construit ce modèle en 1926 pour répondre à son gendre, le
biologiste Umberto D'Ancona, qui demandait pourquoi les poissons prédateurs avaient prospéré dans
l'Adriatique quand la pêche s'était interrompue pendant la Première Guerre mondiale ; Alfred
Lotka avait trouvé les mêmes équations en 1920 par la chimie. La partie (d) est la véritable
réponse à D'Ancona --- ce sont les moyennes, non les instantanés, que le modèle prédit --- et les
orbites fermées font de ce système le système conservatif canonique de la biologie
mathématique.

\begin{refs}
\item Contexte : Équations de prédation de Lotka-Volterra --- Wikipédia (\url{https://fr.wikipedia.org/wiki/%C3%89quations_de_pr%C3%A9dation_de_Lotka-Volterra})
\item Lecture : M. W. Hirsch, S. Smale \& R. L. Devaney, \textit{Differential Equations, Dynamical Systems, and an Introduction to Chaos}, chap. 11
\item Lecture : S. H. Strogatz, \textit{Nonlinear Dynamics and Chaos}, chap. 6
\end{refs}

\bigskip

\begin{problem}[{La période du pendule --- Licence}]
\textbf{ODE-12.} \textbf{Problème.} Le pendule non amorti : $ \theta'' + \frac{g}{L} \sin\theta = 0 $, lâché
sans vitesse initiale à l'amplitude $ \theta_0 \in (0, \pi) $.
\begin{enumerate}
\item[(a)] Démontrer la conservation de l'énergie :
$ E = \frac{1}{2} (\theta')^2 - \frac{g}{L} \cos\theta $ est constante le long des
solutions.
\item[(b)] S'en servir pour établir la période exacte
\[ T(\theta_0) = 4 \sqrt{\frac{L}{g}} \int_0^{\pi/2}
\frac{d\phi}{\sqrt{1 - k^2 \sin^2\phi}}, \qquad k = \sin\frac{\theta_0}{2} , \]
une intégrale elliptique complète de première espèce.
\item[(c)] Démontrer que $ T(\theta_0) $ est strictement croissante, que
$ T \to 2\pi \sqrt{L/g} $ quand $ \theta_0 \to 0 $, et que $ T \to \infty $ quand
$ \theta_0 \to \pi $. Interpréter cette dernière limite dans le portrait de phase (la
séparatrice).
\item[(d)] Montrer que $ T(\theta_0) \approx 2\pi \sqrt{L/g} \left( 1 + \frac{\theta_0^2}{16} \right) $
aux petites amplitudes, ce qui quantifie l'erreur de l'isochronisme de Galilée.
\end{enumerate}

\end{problem}

\noindent Galilée croyait la période du pendule indépendante de l'amplitude ; Huygens,
qui construisait les premières horloges à pendule (1656), savait qu'il n'en était rien et a même
inventé une suspension cycloïdale pour annuler le défaut. L'intégrale elliptique de (b) a résisté
à toutes les tentatives d'évaluation élémentaire et a fini par créer une branche nouvelle des
mathématiques --- les fonctions elliptiques de Legendre, Abel et Jacobi. Ce problème est le pont
entre l'oscillateur linéaire (ODE-05) et la dynamique véritablement non linéaire.

\begin{refs}
\item Contexte : Pendule simple --- Wikipédia (\url{https://fr.wikipedia.org/wiki/Pendule_simple})
\item Lecture : V. I. Arnold, \textit{Ordinary Differential Equations}
\item Lecture : S. H. Strogatz, \textit{Nonlinear Dynamics and Chaos}, chap. 6
\end{refs}

\bigskip

\begin{problem}[{L'identité d'Abel et le wronskien --- Licence, short}]
\textbf{ODE-25.} \textbf{Problème.} Soient $ y_1, y_2 $ deux solutions de
$ y'' + p(t)y' + q(t)y = 0 $ sur un intervalle $ I $, avec $ p, q $ continues, et posons
$ W = y_1 y_2' - y_1' y_2 $. Démontrer l'identité d'Abel
\[ W(t) = W(t_0)\exp\!\left( -\int_{t_0}^{t} p(s)\,ds \right), \]
et en déduire que $ W $ est soit identiquement nul, soit partout non nul --- donc que $ y_1 $ et
$ y_2 $ sont linéairement dépendantes exactement lorsque leur wronskien s'annule en un seul
point de $ I $.
\end{problem}

\noindent Niels Henrik Abel a publié cette identité en 1827 ; le déterminant porte le nom
de Józef Hoene-Wroński, dont Thomas Muir a ranimé la notation des décennies plus tard. C'est la
dichotomie qui fait du wronskien un test utilisable, et elle est propre aux \textit{solutions} d'une
même équation linéaire : les fonctions $ t^2 $ et $ t|t| $ ont un wronskien identiquement nul
sur $\mathbb{R}$ et sont pourtant indépendantes, si bien que l'implication générale est fausse.

\begin{refs}
\item Contexte : Wronskien --- Wikipédia (\url{https://fr.wikipedia.org/wiki/Wronskien})
\item Contexte : Identité d'Abel --- Wikipédia (en anglais) (\url{https://en.wikipedia.org/wiki/Abel%27s_identity})
\item Lecture : W. E. Boyce \& R. C. DiPrima, \textit{Elementary Differential Equations}, chap. 3
\end{refs}

\bigskip

\begin{problem}[{Exactitude et facteur intégrant --- Licence, short}]
\textbf{ODE-26.} \textbf{Problème.} Sur un ouvert simplement connexe
$ U \subseteq \mathbb{R}^2 $, démontrer que $ M(x,y)\,dx + N(x,y)\,dy = 0 $ avec
$ M, N \in C^1(U) $ admet un potentiel $ F $ --- une fonction telle que $ F_x = M $ et
$ F_y = N $, de sorte que les courbes solutions sont les lignes de niveau $ F = c $ --- si et
seulement si $ M_y = N_x $ partout sur $ U $. Montrer ensuite que
\[ (3xy + y^2)\,dx + (x^2 + xy)\,dy = 0 \]
n'est pas exacte, que la multiplication par $ \mu(x) = x $ la rend exacte, et la résoudre.
\end{problem}

\noindent Le critère $ M_y = N_x $ est le cas bidimensionnel de l'énoncé selon lequel
une forme fermée sur un domaine simplement connexe est exacte --- le lemme de Poincaré, rencontrant
les équations différentielles un siècle et demi avant de porter ce nom. Euler et Clairaut ont mis
au point le test et l'astuce du facteur intégrant dans les années 1730 et 1740. Toute équation du
premier ordre admet localement un facteur intégrant ; la difficulté honnête, que cet exemple
dissimule, est qu'en trouver un est d'ordinaire plus dur que l'équation elle-même.

\begin{refs}
\item Contexte : Équation différentielle exacte --- Wikipédia (en anglais) (\url{https://en.wikipedia.org/wiki/Exact_differential_equation})
\item Contexte : Facteur intégrant --- Wikipédia (\url{https://fr.wikipedia.org/wiki/Facteur_int%C3%A9grant})
\item Lecture : W. E. Boyce \& R. C. DiPrima, \textit{Elementary Differential Equations}, chap. 2
\end{refs}

\bigskip

\begin{problem}[{Transformation de Laplace : interrupteurs, coups de marteau et fonction de Green --- Licence}]
\textbf{ODE-27.} \textbf{Problème.} Pour $ f $ continue par morceaux sur $ [0,\infty) $ et d'ordre
exponentiel, la transformée de Laplace est
\[ F(s) = \mathcal{L}\{f\}(s) = \int_0^{\infty} e^{-st} f(t)\,dt . \]
On note $ u_a(t) $ l'échelon unité valant $ 0 $ pour $ t <a $ et $ 1 $ pour
$ t \ge a $.
\begin{enumerate}
\item[(a)] Énoncer des hypothèses sous lesquelles l'intégration par parties est légitime et démontrer
$ \mathcal{L}\{f'\}(s) = sF(s) - f(0) $ ; en déduire la règle pour $ f'' $ et s'en servir pour
résoudre $ y'' + 4y = 0 $, $ y(0) = 1 $, $ y'(0) = 0 $, par la seule transformation.
\item[(b)] Démontrer la règle de retard $ \mathcal{L}\{u_a(t)g(t-a)\}(s) = e^{-as}G(s) $ et résoudre
$ y'' + y = u_{\pi}(t) $ avec $ y(0) = y'(0) = 0 $. Démontrer que la solution est continûment
dérivable mais pas deux fois dérivable en $ t = \pi $, et expliquer pourquoi aucune équation à
forçage discontinu ne peut faire mieux.
\item[(c)] Passons à l'impulsion : résoudre $ y'' + y = \delta(t - \pi) $ avec données initiales
nulles, et dire précisément ce que \og{} solution \fg{} signifie. Montrer que votre fonction est continue,
que sa dérivée présente un saut de taille $ 1 $ en $ t = \pi $, et qu'elle est la limite
quand $ \varepsilon \to 0^{+} $ des solutions forcées par $ \varepsilon^{-1} $ sur
$ [\pi, \pi + \varepsilon] $ et $ 0 $ ailleurs.
\item[(d)] Démontrer le théorème de convolution $ \mathcal{L}\{f * g\} = FG $, où
$ (f*g)(t) = \int_0^t f(t - \tau)g(\tau)\,d\tau $. En déduire que la solution de
$ y'' + y = f(t) $ à données nulles est $ \int_0^t \sin(t - \tau) f(\tau)\,d\tau $, vérifier
que la variation des constantes donne la même formule, et identifier $ \sin(t - \tau) $ comme la
réponse impulsionnelle de l'oscillateur.
\end{enumerate}

\end{problem}

\noindent Oliver Heaviside a inventé ce calcul dans les années 1880 pour analyser les
lignes télégraphiques, en traitant $ d/dt $ comme un symbole algébrique ; les rapporteurs de la
Royal Society objectèrent le manque de rigueur, et il répondit qu'il ne refuserait pas son dîner
sous prétexte qu'il ne comprenait pas la digestion. Bromwich et Carson ont fourni la rigueur
manquante à l'aide d'intégrales de contour dans les années 1910 et 1920, et les travaux de
Laplace lui-même sur les fonctions génératrices, en 1785, avaient fourni la transformation. La
partie (c) est un petit scandale rendu respectable : le delta de Dirac de 1927 fut une fiction
jusqu'à la théorie des distributions de Laurent Schwartz (1945--1950), et la fonction que vous
calculez est la fonction de Green que tout ingénieur appelle fonction de transfert.

\begin{refs}
\item Contexte : Transformation de Laplace --- Wikipédia (\url{https://fr.wikipedia.org/wiki/Transformation_de_Laplace})
\item Contexte : Distribution de Dirac --- Wikipédia (\url{https://fr.wikipedia.org/wiki/Distribution_de_Dirac})
\item Lecture : W. E. Boyce \& R. C. DiPrima, \textit{Elementary Differential Equations}, chap. 6
\item Cours : MIT OCW 18.03SC Differential Equations (\url{https://ocw.mit.edu/courses/18-03sc-differential-equations-fall-2011/})
\end{refs}

\bigskip

\begin{problem}[{Frobenius en un point singulier régulier --- Licence}]
\textbf{ODE-36.} \textbf{Problème.} On considère
\[ x^2 y'' + x\,p(x)\,y' + q(x)\,y = 0 \]
sur $ 0 < x < \rho $, où $ p $ et $ q $ sont données par des séries entières
convergentes sur $ |x| < \rho $ --- le modèle d'un \textit{point singulier régulier} en
$ x = 0 $. L'équation indicielle associée est $ r(r-1) + p(0)\,r + q(0) = 0 $.
\begin{enumerate}
\item[(a)] S'échauffer avec l'équation d'Euler $ x^2 y'' + \alpha x y' + \beta y = 0 $ (constantes
$ \alpha, \beta $). Montrer que $ y = x^r $ en est solution exactement lorsque $ r $
vérifie l'équation indicielle, et établir --- sans la deviner --- la seconde solution
$ x^r \ln x $ lorsque les racines indicielles coïncident.
\item[(b)] Supposer les racines indicielles $ r_1, r_2 $ réelles avec $ r_1 - r_2 $ non entier.
Substituer $ y = x^{r}\sum_{n \ge 0} a_n x^n $ avec $ a_0 = 1 $, établir la récurrence sur
les $ a_n $, et démontrer qu'elle les détermine de manière unique pour $ r = r_1 $ et pour
$ r = r_2 $. Montrer que les deux solutions obtenues sont linéairement indépendantes, et
énoncer (sans le démontrer) le théorème garantissant la convergence des séries pour
$ |x| < \rho $.
\item[(c)] Voici maintenant le cas qui a fait la célébrité de la méthode. Pour l'équation de Bessel
d'ordre zéro,
\[ x^2 y'' + x y' + x^2 y = 0, \]
montrer que l'équation indicielle est $ r^2 = 0 $, calculer la solution de Frobenius
\[ J_0(x) = \sum_{k \ge 0} \frac{(-1)^k}{(k!)^2}\left(\frac{x}{2}\right)^{2k}, \]
et démontrer qu'il n'existe pas de seconde solution de la forme pure $ x^{r}\sum a_n x^n $.
Montrer ensuite qu'une seconde solution doit contenir un terme $ J_0(x)\ln x $, en substituant
$ y = J_0(x)\ln x + \sum_{n \ge 1} b_n x^n $ et en trouvant une récurrence pour les
$ b_n $.
\item[(d)] Montrer que l'hypothèse compte. Vérifier que $ y' = y/x^2 $ admet les solutions
$ y = C e^{-1/x} $ ; démontrer qu'aucune solution non nulle ne peut s'écrire comme $ x^r $
multiplié par une série entière convergente au voisinage de $ 0 $, et identifier le trait de
l'équation qui la place hors de la théorie des parties (a)--(c).
\end{enumerate}

\end{problem}

\noindent Lazarus Fuchs a identifié en 1866 les équations dont les singularités sont
assez douces pour que les séries fonctionnent --- les points singuliers réguliers --- et Georg
Frobenius, dans \og{}~Über die Integration der linearen Differentialgleichungen durch
Reihen~\fg{} (\textit{Journal de Crelle} 76, 1873, 214--235), en a donné la construction
systématique, logarithmes compris. Le bénéfice, c'est toute la troupe des fonctions spéciales du
XIXe siècle : les fonctions de Bessel pour le tambour circulaire et le cylindre,
celles de Legendre pour la sphère (ODE-37), et la série hypergéométrique de Gauss, chacune une
série de Frobenius en un point singulier de son équation. La leçon de (d) est de savoir où la
théorie s'arrête : en un point singulier irrégulier, les solutions peuvent avoir des singularités
essentielles, et les développements naturels y sont des séries asymptotiques divergentes plutôt
que convergentes.

\begin{refs}
\item Contexte : Méthode de Frobenius --- Wikipédia (\url{https://fr.wikipedia.org/wiki/M%C3%A9thode_de_Frobenius})
\item Contexte : Point singulier régulier --- Wikipédia (\url{https://fr.wikipedia.org/wiki/Point_singulier_r%C3%A9gulier})
\item Lecture : W. E. Boyce \& R. C. DiPrima, \textit{Elementary Differential Equations}, chap. 5
\item Lecture : E. L. Ince, \textit{Ordinary Differential Equations} (Dover)
\end{refs}

\bigskip

\begin{problem}[{L'équation de Legendre --- Licence, short}]
\textbf{ODE-37.} \textbf{Problème.} Pour l'équation de Legendre
\[ (1 - x^2)\,y'' - 2x\,y' + \lambda\,y = 0, \qquad -1 < x < 1, \]
développer $ y $ en série entière autour de $ x = 0 $, établir la récurrence à deux termes sur
les coefficients, et démontrer qu'une solution bornée sur tout $ (-1, 1) $ et de limites finies
en $ x = \pm 1 $ existe exactement lorsque $ \lambda = n(n+1) $ pour un entier $ n \ge 0 $,
la série se terminant alors en un polynôme $ P_n $ de degré $ n $. Écrire ensuite l'équation
sous la forme auto-adjointe $ \big( (1-x^2)\,y' \big)' + \lambda\,y = 0 $ et s'en servir pour
démontrer que $ \int_{-1}^{1} P_n(x) P_m(x)\,dx = 0 $ dès que $ n \ne m $.
\end{problem}

\noindent Adrien-Marie Legendre a introduit ces polynômes dans les années 1780 en
calculant l'attraction gravitationnelle des sphéroïdes : ce sont les coefficients du développement
du potentiel newtonien $ 1/|x - y| $ en puissances du rapport des distances, si bien qu'ils sont
inscrits dans tout développement multipolaire de la physique. Olinde Rodrigues a donné la formule
close $ P_n = \frac{1}{2^n n!}\frac{d^n}{dx^n}(x^2-1)^n $ en 1816. Le schéma de ce problème --- la
bornitude aux extrémités singulières sélectionne un ensemble discret de $ \lambda $, et les
fonctions propres obtenues sont orthogonales --- est le prototype de la théorie de Sturm--Liouville,
et c'est ce qui fait de ces polynômes la bonne base sur la sphère.

\begin{refs}
\item Contexte : Polynôme de Legendre --- Wikipédia (\url{https://fr.wikipedia.org/wiki/Polyn%C3%B4me_de_Legendre})
\item Contexte : Formule de Rodrigues --- Wikipédia (\url{https://fr.wikipedia.org/wiki/Formule_de_Rodrigues})
\item Lecture : W. E. Boyce \& R. C. DiPrima, \textit{Elementary Differential Equations}, chap. 5
\end{refs}

\bigskip

\begin{problem}[{La raideur, et pourquoi les méthodes implicites existent --- Licence}]
\textbf{ODE-38.} \textbf{Problème.} Appliquer des méthodes numériques à l'équation test de Dahlquist
$ y' = \lambda y $ avec $ \operatorname{Re}\lambda < 0 $ et un pas $ h > 0 $, de sorte
que la solution exacte décroisse vers $ 0 $.
\begin{enumerate}
\item[(a)] Euler explicite (progressif) donne $ y_{n+1} = (1 + h\lambda)\,y_n $. Démontrer que
$ y_n \to 0 $ pour toute valeur initiale si et seulement si $ |1 + h\lambda| < 1 $, et en
déduire que pour $ \lambda < 0 $ réel la méthode est inutilisable à moins que
$ h < 2/|\lambda| $ --- une restriction imposée par la \textit{stabilité}, non par la
précision.
\item[(b)] Euler implicite (rétrograde) donne $ y_{n+1} = y_n/(1 - h\lambda) $. Démontrer que
$ y_n \to 0 $ pour tout $ h > 0 $ dès que $ \operatorname{Re}\lambda < 0 $ : la
méthode est \textit{A-stable}. Dire précisément quel travail supplémentaire coûte désormais chaque
pas pour un système non linéaire.
\item[(c)] Rendre la douleur concrète. Soit
\[ A = \frac{1}{2}\begin{pmatrix} -1001 & 999 \\ 999 & -1001 \end{pmatrix}, \]
et vérifier que $ A $ a pour valeurs propres $ -1 $ et $ -1000 $, de vecteurs propres
$ (1,1) $ et $ (1,-1) $. Écrire la solution exacte de $ x' = Ax $, montrer qu'après
$ t \approx 0{,}01 $ elle est indiscernable d'une courbe lisse variant sur l'échelle de temps
$ 1 $, et démontrer qu'Euler explicite exige néanmoins $ h < 0{,}002 $ pour toute
l'intégration. Compter les pas nécessaires sur $ [0, 10] $ pour chacune des deux méthodes, et
expliquer où passe le travail supplémentaire dans chaque cas.
\item[(d)] Définir le domaine de stabilité absolue d'une méthode à un pas comme l'ensemble des
$ z = h\lambda $ complexes pour lesquels la solution numérique de l'équation test reste bornée,
et dire la méthode A-stable si ce domaine contient tout le demi-plan gauche. Démontrer qu'aucune
méthode de Runge--Kutta explicite n'est A-stable : montrer que sa fonction de stabilité
$ R(z) $ est un polynôme en $ z $, de sorte que $ |R(z)| \to \infty $ quand
$ |z| \to \infty $. Énoncer ensuite la seconde barrière de Dahlquist pour les méthodes
linéaires à pas multiples (1963) et expliquer, en un paragraphe, pourquoi les méthodes implicites
ne sont pas une inefficacité qu'on pourrait faire disparaître par l'ingénierie mais une nécessité
mathématique.
\end{enumerate}

\end{problem}

\noindent Charles Curtiss et Joseph Hirschfelder ont forgé le mot \og{}~stiff~\fg{}
(raide) en 1952, dans une note aux \textit{Proceedings of the National Academy of Sciences}, après
que des cinétiques chimiques aux constantes de vitesse follement séparées eurent mis en échec les
intégrateurs standard de l'époque. Germund Dahlquist a rendu le phénomène exact en 1963 en
définissant l'A-stabilité et en démontrant qu'elle se paie : une méthode linéaire à pas multiples
A-stable est d'ordre au plus deux, la règle des trapèzes étant la meilleure d'entre elles. Tous
les solveurs raides en usage --- les formules de différentiation rétrograde de Gear, LSODE,
l'\textit{ode15s} de MATLAB --- descendent de ce marché, et la morale mérite d'être énoncée clairement :
un problème raide n'est pas un problème difficile, c'est un problème dont l'échelle de temps
intéressante et l'échelle de temps limitant la stabilité diffèrent de plusieurs ordres de
grandeur. La même structure à deux échelles, vue analytiquement plutôt que numériquement, est la
couche limite d'ODE-41.

\begin{refs}
\item Contexte : Équation différentielle raide --- Wikipédia (\url{https://fr.wikipedia.org/wiki/%C3%89quation_diff%C3%A9rentielle_raide})
\item Contexte : Méthode d'Euler implicite --- Wikipédia (en anglais) (\url{https://en.wikipedia.org/wiki/Backward_Euler_method})
\item Article : C. F. Curtiss \& J. O. Hirschfelder, \og{} Integration of stiff equations \fg{}, \textit{Proceedings of the National Academy of Sciences} 38 (1952), 235--243
\item Lecture : E. Hairer \& G. Wanner, \textit{Solving Ordinary Differential Equations II: Stiff and Differential-Algebraic Problems}
\end{refs}

\bigskip


\section*{Master}

\begin{problem}[{Fonctions de Liapounov --- Master}]
\textbf{ODE-13.} \textbf{Problème.} Soit $ x^* $ un point d'équilibre de $ x' = f(x) $ dans
$ \mathbb{R}^n $, et soit $ V $ une fonction lisse telle que $ V(x^*) = 0 $ et
$ V >0 $ au voisinage (définie positive). On note $ \dot{V} $ la dérivée de $ V $ le
long des solutions.
\begin{enumerate}
\item[(a)] Démontrer les théorèmes de Liapounov : si $ \dot{V} \le 0 $ près de $ x^* $, l'équilibre
est stable ; si de plus $ \dot{V} <0 $ sauf en $ x^* $, il est asymptotiquement
stable.
\item[(b)] Appliquer au pendule amorti $ \theta'' + b\theta' + \sin\theta = 0 $ avec l'énergie
$ V = \frac{1}{2}(\theta')^2 + 1 - \cos\theta $ : montrer que $ \dot{V} \le 0 $ mais
\textit{pas} strictement négative, si bien que (a) seule ne donne que la stabilité.
\item[(c)] Énoncer le principe d'invariance de LaSalle et s'en servir pour conclure malgré tout à la
stabilité asymptotique de l'équilibre bas.
\item[(d)] Expliquer le marché : la méthode de Liapounov n'exige aucune formule de solution, mais il
n'existe aucune recette générale pour trouver $ V $. Construire $ V $ pour le système
gradient $ x' = -\nabla U(x) $ et démontrer que toute solution bornée converge vers l'ensemble
des points critiques de $ U $.
\end{enumerate}

\end{problem}

\noindent La thèse de doctorat d'Alexandre Liapounov, en 1892, \og{} Le problème général de
la stabilité du mouvement \fg{}, a fait de la théorie de la stabilité une mathématique ; sa
\og{} seconde méthode \fg{} --- démontrer la stabilité en exhibant une énergie généralisée qui ne peut
croître --- est aujourd'hui partout, de l'automatique (tout certificat de pilote automatique est une
fonction de Liapounov) à la biologie mathématique. Le raffinement de LaSalle (1960) traite
exactement la dégénérescence que vous rencontrez en (b).

\begin{refs}
\item Contexte : Fonction de Liapounov --- Wikipédia (\url{https://fr.wikipedia.org/wiki/Fonction_de_Liapounov})
\item Lecture : M. W. Hirsch, S. Smale \& R. L. Devaney, \textit{Differential Equations, Dynamical Systems, and an Introduction to Chaos}, chap. 9
\item Lecture : S. H. Strogatz, \textit{Nonlinear Dynamics and Chaos}, chap. 7
\end{refs}

\bigskip

\begin{problem}[{Piéger un cycle limite : Poincaré--Bendixson --- Master}]
\textbf{ODE-14.} \textbf{Problème.}
\begin{enumerate}
\item[(a)] Énoncer le théorème de Poincaré--Bendixson : une orbite positive d'un système plan de classe
$ C^1 $ qui reste dans un compact sans point d'équilibre possède une orbite périodique dans son
ensemble limite.
\item[(b)] L'appliquer à $ x' = x - y - x^3 $, $ y' = x + y - y^3 $ : démontrer que l'origine est le
seul point d'équilibre et qu'il est répulsif ; montrer, en estimant la dérivée de
$ r^2 = x^2 + y^2 $ le long des solutions, qu'un anneau convenable centré à l'origine est
positivement invariant ; et conclure que le système possède au moins une orbite périodique.
\item[(c)] Démontrer le critère négatif de Bendixson : si la divergence du champ de vecteurs garde un
signe strict sur une région simplement connexe, cette région ne contient aucune orbite
périodique. (Théorème de Green.) S'en servir pour montrer que le pendule amorti (ODE-13b) n'a
aucune orbite périodique dans toute bande où il s'applique.
\item[(d)] Expliquer pourquoi la méthode tout entière est bidimensionnelle : quel fait topologique
concernant le plan la démonstration de (a) utilise-t-elle, et à quoi son échec dans
$ \mathbb{R}^3 $ laisse-t-il la place ? (ODE-17 répond.)
\end{enumerate}

\end{problem}

\noindent Poincaré a créé la théorie des cycles limites dans les années 1880 ; Ivar
Bendixson en a fourni la démonstration topologique épurée en 1901. Le théorème est la dichotomie
fondamentale de la dynamique plane --- les orbites bornées se rangent vers des équilibres ou des
cycles, rien de plus étrange --- et son échec en dimension trois est précisément là où vit le chaos.
Le théorème de Jordan qui travaille silencieusement en (d) est une leçon sur la façon dont la
topologie contraint la dynamique.

\begin{refs}
\item Contexte : Théorème de Poincaré-Bendixson --- Wikipédia (\url{https://fr.wikipedia.org/wiki/Th%C3%A9or%C3%A8me_de_Poincar%C3%A9-Bendixson})
\item Lecture : S. H. Strogatz, \textit{Nonlinear Dynamics and Chaos}, chap. 7
\item Lecture : M. W. Hirsch, S. Smale \& R. L. Devaney, \textit{Differential Equations, Dynamical Systems, and an Introduction to Chaos}, chap. 10
\end{refs}

\bigskip

\begin{problem}[{L'oscillateur de Van der Pol --- Master}]
\textbf{ODE-15.} \textbf{Problème.} L'équation de Van der Pol, pour $ \mu >0 $ :
\[ x'' - \mu (1 - x^2)\, x' + x = 0 . \]
(a) Interpréter la structure des signes : l'amortissement est négatif pour $ |x| <1 $
(énergie injectée) et positif pour $ |x| >1 $ (énergie retirée). Montrer que l'origine est
l'unique point d'équilibre et qu'il est instable pour tout $ \mu >0 $.
(b) Énoncer le théorème de Liénard sur les équations $ x'' + f(x) x' + g(x) = 0 $, en vérifier
les hypothèses pour Van der Pol, et conclure : l'équation possède exactement une orbite
périodique, et toute autre trajectoire non stationnaire s'enroule vers elle --- un cycle limite
unique et globalement attractif.
(c) Problème de compréhension : pour $ \mu \gg 1 $, le cycle devient une \textit{oscillation de
relaxation} --- de longues reptations lentes ponctuées de sauts quasi instantanés. Analyser la
structure lent--rapide dans le plan de Liénard et en déduire la période au premier ordre
$ T \sim (3 - 2\ln 2)\, \mu $.
\end{problem}

\noindent Balthasar van der Pol a trouvé cette équation dans les années 1920 dans les
circuits à lampes chez Philips ; avec van der Mark il en a tiré un modèle électrique du battement
cardiaque (1928), et leur note sur le \og{} bruit irrégulier \fg{} entendu à certaines fréquences
d'excitation est l'une des premières observations expérimentales du chaos jamais consignées.
L'équation est l'oscillateur auto-entretenu prototype --- le cœur mathématique des stimulateurs
cardiaques, des cordes de violon et de la dynamique des lasers --- et l'application standard du
théorème d'unicité de Liénard (1928).

\begin{refs}
\item Contexte : Oscillateur de Van der Pol --- Wikipédia (\url{https://fr.wikipedia.org/wiki/Oscillateur_de_Van_der_Pol})
\item Lecture : S. H. Strogatz, \textit{Nonlinear Dynamics and Chaos}, chap. 7
\item Article : B. van der Pol, \og{} On relaxation-oscillations \fg{}, \textit{Philosophical Magazine} (7) 2 (1926), 978--992
\end{refs}

\bigskip

\begin{problem}[{La suite logistique et la route vers le chaos par doublement de période --- Master}]
\textbf{ODE-16.} \textbf{Problème.} On itère $ x_{n+1} = r x_n (1 - x_n) $ sur
$ [0, 1] $, $ 0 \le r \le 4 $.
\begin{enumerate}
\item[(a)] Trouver les points fixes et démontrer : $ x = 0 $ est attractif pour $ r <1 $ ; le
point fixe $ 1 - 1/r $ est attractif pour $ 1 <r <3 $ (utiliser le critère de la
dérivée, et démontrer ce critère).
\item[(b)] Démontrer qu'en $ r = 3 $ le point fixe perd sa stabilité et qu'une orbite de période 2
naît : calculer explicitement les points de période 2 à partir d'une équation du second degré,
et montrer qu'ils sont attractifs pour $ 3 <r <1 + \sqrt{6} $.
\item[(c)] Problème de compréhension : décrire la cascade de doublements de période --- les paramètres
$ r_1 <r_2 <\cdots \to r_\infty \approx 3{,}5699 $ où apparaissent les orbites de période
$ 2^k $ --- et expliquer le sens de la constante de Feigenbaum
$ \delta = \lim_k (r_k - r_{k-1})/(r_{k+1} - r_k) \approx 4{,}6692 $ : le même nombre apparaît
pour toute une classe d'applications (universalité). Expliquer, au niveau des idées, comment la
renormalisation en rend compte, et noter que les conjectures de Feigenbaum ont été démontrées par
Lanford (1982) au moyen d'une démonstration assistée par ordinateur.
\item[(d)] Démontrer qu'en $ r = 4 $ l'application est semi-conjuguée au doublement
$ \theta \mapsto 2\theta \bmod 2\pi $ via $ x = \sin^2(\theta/2) $, et en déduire la
sensibilité aux conditions initiales.
\end{enumerate}

\end{problem}

\noindent La revue de Robert May dans \textit{Nature} en 1976 a rendu célèbre cette
équation d'une seule ligne : le modèle le plus simple de l'écologie contient une cascade infinie
de doublements de période, puis le chaos. Mitchell Feigenbaum a découvert la constante universelle
4,6692\dots{} en 1975 avec une calculatrice de poche, et l'universalité --- des systèmes radicalement
différents partageant un comportement quantitatif exact --- a été bientôt confirmée dans de vraies
expériences de mécanique des fluides (Libchaber). La démonstration par Oscar Lanford, en 1982, des
conjectures de Feigenbaum fut un jalon précoce de l'analyse rigoureuse assistée par ordinateur,
préfigurant ODE-17.

\begin{refs}
\item Contexte : Suite logistique --- Wikipédia (\url{https://fr.wikipedia.org/wiki/Suite_logistique})
\item Contexte : Nombres de Feigenbaum --- Wikipédia (\url{https://fr.wikipedia.org/wiki/Nombres_de_Feigenbaum})
\item Lecture : S. H. Strogatz, \textit{Nonlinear Dynamics and Chaos}, chap. 10
\item Article : R. M. May, \og{} Simple mathematical models with very complicated dynamics \fg{}, \textit{Nature} 261 (1976), 459--467
\end{refs}

\bigskip

\begin{problem}[{Le système de Lorenz : ce que signifie \og{} l'attracteur existe \fg{} --- Master}]
\textbf{ODE-17.} \textbf{Problème.} Le système de Lorenz, avec les paramètres classiques
$ \sigma = 10 $, $ \rho = 28 $, $ \beta = 8/3 $ :
\[ x' = \sigma (y - x), \qquad y' = x (\rho - z) - y, \qquad z' = x y - \beta z . \]
(a) Démontrer que le flot contracte les volumes au taux exponentiel constant
$ e^{-(\sigma + 1 + \beta) t} $ (calculer la divergence), de sorte que l'ensemble attracteur est
de volume nul.
(b) Démontrer qu'il existe une région piège : montrer que
$ V = \rho x^2 + \sigma y^2 + \sigma (z - 2\rho)^2 $ vérifie $ \dot{V} <0 $ en dehors d'un
ellipsoïde borné, de sorte qu'une grande boule est atteinte par toute solution et jamais
quittée.
(c) Trouver les trois points d'équilibre pour ces paramètres et montrer qu'ils sont tous
instables.
(d) Problème de compréhension : les parties (a)--(c) disent que les orbites sont bornées sans rien
de stable où se poser --- mais elles ne démontrent \textit{pas} l'existence d'un attracteur étrange.
Expliquer ce que \og{} l'attracteur de Lorenz existe \fg{} affirme de plus (un attracteur robuste portant
une dynamique chaotique, l'hyperbolicité singulière, le modèle géométrique de Lorenz), pourquoi
Smale en a fait son 14e problème pour le XXIe siècle, et comment la
démonstration assistée par ordinateur de Warwick Tucker en 2002 --- formes normales près de
l'origine et intégration rigoureuse par arithmétique d'intervalles --- a tranché la question.
Qu'est-ce qui rend une telle démonstration rigoureuse là où une simulation en virgule flottante ne
l'est pas ?
\end{problem}

\noindent Edward Lorenz, météorologue au MIT, a découvert en 1963 qu'un modèle grossier
de convection à trois modes refusait d'être prévisible : l'effet papillon. Pendant près de
quarante ans, tout le monde pouvait \textit{voir} l'attracteur à l'écran et personne ne pouvait
démontrer qu'il était vraiment là --- les simulations auraient pu, en principe, suivre l'ombre d'un
très long cycle stable. La démonstration de Tucker a comblé cette lacune et fait figure de cas
d'école de l'usage des ordinateurs en mathématiques rigoureuses. Le chaos déterministe, entrevu
d'abord par Poincaré (ODE-18), devient ici un objet tridimensionnel concret.

\begin{refs}
\item Contexte : Attracteur de Lorenz --- Wikipédia (\url{https://fr.wikipedia.org/wiki/Attracteur_de_Lorenz})
\item Lecture : S. H. Strogatz, \textit{Nonlinear Dynamics and Chaos}, chap. 9
\item Article : E. N. Lorenz, \og{} Deterministic nonperiodic flow \fg{}, \textit{Journal of the Atmospheric Sciences} 20 (1963), 130--141
\item Article : W. Tucker, \og{} A rigorous ODE solver and Smale's 14th problem \fg{}, \textit{Foundations of Computational Mathematics} 2 (2002), 53--117
\end{refs}

\bigskip

\begin{problem}[{Le problème des trois corps --- Master}]
\textbf{ODE-18.} \textbf{Problème.} Trois masses ponctuelles s'attirent mutuellement par la gravitation
newtonienne.
\begin{enumerate}
\item[(a)] Un îlot soluble : démontrer la solution équilatérale de Lagrange pour des masses égales ---
trois masses égales aux sommets d'un triangle équilatéral, chacune parcourant un cercle autour du
centre de masse commun, satisfont les équations du mouvement --- et déterminer la vitesse angulaire
$ \omega $ requise en fonction des masses et de la longueur du côté.
\item[(b)] Problème de compréhension : lire l'histoire du mémoire de Poincaré pour le prix du roi
Oscar II (1889) --- y compris l'erreur qu'il a trouvée après l'attribution du prix, et la découverte
des enchevêtrements homoclines dans le mémoire corrigé --- et expliquer avec vos propres mots ce que
Poincaré a vu et pourquoi cela a tué l'espoir de résoudre la mécanique par des formules.
\item[(c)] Problème de compréhension : Karl Sundman (1912) a néanmoins produit une série entière
convergente en $ t^{1/3} $ résolvant le problème des trois corps pour toutes les conditions
initiales de moment cinétique non nul. Concilier (b) et (c) : en quel sens celle de Sundman est-elle
une \og{}~solution complète~\fg{}, et pourquoi est-elle inutilisable en pratique
(convergence si lente qu'un usage astronomique exigerait un nombre inimaginable de termes) et
muette sur les questions qualitatives ?
\end{enumerate}

\end{problem}

\noindent Le problème des trois corps est là où le rêve d'une mécanique céleste exacte est
mort et où les systèmes dynamiques sont nés. Le mémoire corrigé de Poincaré (1890) et ses
\textit{Méthodes nouvelles} en trois volumes (1892--1899) contiennent le premier chaos dont
l'existence ait jamais été démontrée ; l'histoire de June Barrow-Green reconstitue l'épisode
entier, erreur comprise. La série de Sundman, véritable triomphe mathématique, est l'exemple
standard montrant qu'\og{} il existe une solution en série convergente \fg{} et \og{} nous comprenons la
dynamique \fg{} sont deux affirmations différentes --- la motivation même de la théorie qualitative que
cette page enseigne.

\begin{refs}
\item Contexte : Problème des trois corps --- Wikipédia (en anglais) (\url{https://en.wikipedia.org/wiki/Three-body_problem})
\item Lecture : J. Barrow-Green, \textit{Poincaré and the Three Body Problem} (AMS/LMS, 1997)
\item Lecture : V. I. Arnold, \textit{Ordinary Differential Equations}
\end{refs}

\bigskip

\begin{problem}[{Un critère qui interdit les cycles : Bendixson--Dulac --- Master, short}]
\textbf{ODE-28.} \textbf{Problème.} Soit $ D \subseteq \mathbb{R}^2 $ un ouvert simplement
connexe et soient $ P, Q \in C^1(D) $. Démontrer le critère de Bendixson--Dulac : s'il existe
$ \varphi \in C^1(D) $ avec $ \varphi >0 $ telle que
\[ \frac{\partial(\varphi P)}{\partial x} + \frac{\partial(\varphi Q)}{\partial y} \]
garde un signe strict sur tout $ D $, alors $ x' = P(x,y) $, $ y' = Q(x,y) $ n'a aucune
orbite périodique contenue dans $ D $ ; s'en servir ensuite, avec $ \varphi \equiv 1 $, pour
démontrer que le pendule amorti $ x' = y $, $ y' = -\sin x - cy $ avec $ c >0 $ n'a
aucune orbite périodique.
\end{problem}

\noindent Ivar Bendixson a démontré le cas $ \varphi \equiv 1 $ en 1901 et Henri Dulac
a fourni la fonction de poids ; la démonstration est le théorème de Green appliqué à la région
délimitée par un cycle hypothétique, ce qui explique très exactement pourquoi la simple connexité
n'est pas un ornement. C'est le compagnon négatif du théorème de Poincaré--Bendixson d'ODE-14 :
l'un piège les cycles, l'autre les interdit. Trouver une fonction de Dulac est un art sans
algorithme, et la difficulté de compter les cycles en l'absence d'une telle fonction est la porte
d'entrée d'ODE-19.

\begin{refs}
\item Contexte : Théorème de Bendixson--Dulac --- Wikipédia (en anglais) (\url{https://en.wikipedia.org/wiki/Bendixson%E2%80%93Dulac_theorem})
\item Lecture : S. H. Strogatz, \textit{Nonlinear Dynamics and Chaos}, chap. 7
\item Lecture : L. Perko, \textit{Differential Equations and Dynamical Systems}, chap. 3
\end{refs}

\bigskip

\begin{problem}[{Comparaison de Sturm et zéros de la fonction d'Airy --- Master, short}]
\textbf{ODE-29.} \textbf{Problème.} Soient $ q_1 \le q_2 $ continues sur un intervalle
$ I $, soit $ u \not\equiv 0 $ solution de $ u'' + q_1 u = 0 $ et soit
$ v \not\equiv 0 $ solution de $ v'' + q_2 v = 0 $. Démontrer le théorème de comparaison de
Sturm --- entre deux zéros consécutifs de $ u $ se trouve un zéro de $ v $, sauf si
$ q_1 \equiv q_2 $ là et si $ v $ est un multiple constant de $ u $ --- et en déduire que
toute solution non triviale de l'équation d'Airy $ y'' = xy $ possède une infinité de zéros sur
$ x <0 $ et au plus un zéro sur $ x >0 $.
\end{problem}

\noindent Le mémoire de 1836 de Charles-François Sturm sur les équations linéaires du
second ordre est habituellement tenu pour la première théorie purement qualitative des équations
différentielles : il a extrait l'oscillation des solutions du signe du coefficient sans rien
résoudre. L'équation d'Airy, introduite par George Biddell Airy en 1838 pour expliquer
l'intensité lumineuse près d'une caustique --- les arcs surnuméraires de l'arc-en-ciel --- est le
modèle standard d'un \textit{point de retournement}, où une solution passe de l'oscillation à la
croissance, et le théorème de Sturm localise exactement ce basculement.

\begin{refs}
\item Contexte : Théorème de comparaison de Sturm-Picone --- Wikipédia (\url{https://fr.wikipedia.org/wiki/Th%C3%A9or%C3%A8me_de_comparaison_de_Sturm-Picone})
\item Contexte : Fonction d'Airy --- Wikipédia (\url{https://fr.wikipedia.org/wiki/Fonction_d%27Airy})
\item Lecture : E. A. Coddington \& N. Levinson, \textit{Theory of Ordinary Differential Equations}
\end{refs}

\bigskip

\begin{problem}[{Systèmes hamiltoniens : ce que la conservation interdit --- Master}]
\textbf{ODE-30.} \textbf{Problème.} Un système hamiltonien sur $ \mathbb{R}^{2n} $ de hamiltonien lisse
$ H(q,p) $ s'écrit
\[ q_i' = \frac{\partial H}{\partial p_i}, \qquad p_i' = -\frac{\partial H}{\partial q_i},
\qquad i = 1, \dots, n. \]
\begin{enumerate}
\item[(a)] Démontrer que $ H $ est constant le long de toute trajectoire, et démontrer le théorème de
Liouville : le champ de vecteurs est de divergence nulle, et par conséquent le flot au temps
$ t $ préserve la mesure de Lebesgue sur $ \mathbb{R}^{2n} $. (Dériver le déterminant
jacobien du flot.)
\item[(b)] En déduire deux interdictions. Démontrer qu'un système hamiltonien n'a aucun point
d'équilibre asymptotiquement stable ni aucune orbite périodique attractive --- donc aucun cycle
limite au sens d'ODE-14 --- et conclure que le pendule amorti d'ODE-28 n'admet aucune formulation
hamiltonienne dans ces coordonnées.
\item[(c)] Démontrer le théorème de récurrence de Poincaré dans ce cadre : si un ensemble mesurable
$ A $ de mesure finie est invariant par le flot, alors presque tout point de $ A $ revient
dans tout voisinage de lui-même à des temps arbitrairement grands. Dire exactement quelles
hypothèses --- préservation de la mesure, finitude --- chaque étape consomme.
\item[(d)] Rendre cela concret avec le pendule non amorti $ H(q,p) = \tfrac{1}{2}p^2 - \cos q $.
Esquisser les lignes de niveau, identifier la séparatrice passant par $ (q,p) = (\pi, 0) $, et
montrer que les orbites périodiques se présentent en familles continues à un paramètre --- aucune
n'est donc isolée, en accord avec (b).
\end{enumerate}

\end{problem}

\noindent William Rowan Hamilton a refondu la mécanique sous cette forme en 1834--1835 ;
Liouville a démontré le théorème sur les volumes en 1838, et le théorème de récurrence de Poincaré
est paru dans son mémoire primé de 1890 sur le problème des trois corps (ODE-18). La récurrence
est aussitôt devenue célèbre : Zermelo l'a retournée contre Boltzmann en 1896, arguant qu'un gaz
obéissant à la mécanique doit finir par revenir à son état initial et ne peut donc obéir à un
théorème $ H $ irréversible --- les temps de récurrence sont simplement absurdes, non infinis. La
partie (b) explique pourquoi frottement et conservation sont mathématiquement incompatibles, et le
théorème de Noether de 1918 complète le tableau en déduisant chaque grandeur conservée d'une
symétrie.

\begin{refs}
\item Contexte : Mécanique hamiltonienne --- Wikipédia (\url{https://fr.wikipedia.org/wiki/M%C3%A9canique_hamiltonienne})
\item Contexte : Théorème de Liouville (hamiltonien) --- Wikipédia (\url{https://fr.wikipedia.org/wiki/Th%C3%A9or%C3%A8me_de_Liouville_(hamiltonien)})
\item Contexte : Théorème de récurrence de Poincaré --- Wikipédia (\url{https://fr.wikipedia.org/wiki/Th%C3%A9or%C3%A8me_de_r%C3%A9currence_de_Poincar%C3%A9})
\item Lecture : V. I. Arnold, \textit{Mathematical Methods of Classical Mechanics}, chap. 3 et 8
\end{refs}

\bigskip

\begin{problem}[{Théorie de Floquet et la balançoire qu'on pompe --- Master}]
\textbf{ODE-39.} \textbf{Problème.} Soit $ A(t) $ continue et $ T $-périodique, et soit $ \Phi(t) $ la
matrice fondamentale de $ x' = A(t)x $ avec $ \Phi(0) = I $. La \textit{matrice de monodromie}
est $ M = \Phi(T) $, et ses valeurs propres sont les \textit{multiplicateurs de Floquet}.
\begin{enumerate}
\item[(a)] Démontrer $ \Phi(t + T) = \Phi(t)\,M $ pour tout $ t $, en montrant que les deux membres
résolvent le même problème de Cauchy matriciel. En déduire la représentation de Floquet
$ \Phi(t) = P(t)\,e^{tB} $ avec $ P $ périodique et $ B $ constante (on admettra la
$ 2T $-périodicité si l'on tient à des matrices réelles), et démontrer que la solution nulle est
asymptotiquement stable si tout multiplicateur est dans le disque unité ouvert, et instable si un
multiplicateur est à l'extérieur.
\item[(b)] Démontrer la formule de Liouville
$ \det \Phi(t) = \exp \int_0^t \operatorname{tr} A(s)\,ds $. L'appliquer à l'équation de Hill
$ x'' + q(t)\,x = 0 $ avec $ q $ périodique, écrite comme système plan : montrer que
$ \det M = 1 $, de sorte que les multiplicateurs sont $ \rho $ et $ 1/\rho $, et conclure que
toutes les solutions ne sont bornées que si $ |\operatorname{tr} M| \le 2 $, tandis que
$ |\operatorname{tr} M| > 2 $ force l'existence de solutions à croissance
exponentielle.
\item[(c)] Interpréter la balançoire. Dans l'équation de Mathieu
$ x'' + (\delta + \varepsilon \cos t)\,x = 0 $, c'est la force de rappel qui est modulée
périodiquement, et non une force appliquée directement. Expliquer pourquoi un enfant qui monte et
descend son centre de masse deux fois par balancement module au double de la fréquence propre, et
décrire les langues d'instabilité qui, pour $ \varepsilon $ petit, émanent de
$ \delta = n^2/4 $, $ n = 1, 2, 3, \ldots $
\item[(d)] Démontrer que la résonance paramétrique se produit réellement, dans le seul cas calculable à
la main. Prendre l'équation à créneaux (Meissner) : $ q(t) = \omega^2(1 + \varepsilon) $ sur
$ [0, T/2) $ et $ q(t) = \omega^2(1 - \varepsilon) $ sur $ [T/2, T) $, prolongée par
périodicité. Écrire $ M $ comme le produit des deux matrices de transition à coefficients
constants, calculer $ \operatorname{tr} M $ explicitement, et exhiber des paramètres
$ (\omega, \varepsilon, T) $ pour lesquels $ |\operatorname{tr} M| > 2 $ --- une équation à
coefficients bornés et purement oscillants dont toutes les solutions explosent.
\end{enumerate}

\end{problem}

\noindent Gaston Floquet a publié le théorème de représentation en 1883 ; George Hill
avait déjà, en 1877, rencontré le même problème dans le mouvement du périgée lunaire et inventé
les déterminants infinis pour l'attaquer, et Émile Mathieu était arrivé à son équation en 1868 à
partir des vibrations d'une membrane elliptique. La physique est plus ancienne qu'eux tous :
Faraday a vu en 1831 qu'une surface liquide vibrée se ride à la moitié de la fréquence
d'excitation, et Melde en 1859 a excité une corde par un diapason fixé le long de celle-ci et l'a
vue osciller à la moitié de la fréquence du diapason. La résonance paramétrique est aujourd'hui
un outil de travail --- le pendule inversé de Kapitza, stabilisé la tête en bas par une vibration
verticale (1951), l'amplificateur paramétrique, et le piège à ions de Paul, dont le diagramme de
stabilité est exactement le jeu de langues de la partie (c).

\begin{refs}
\item Contexte : Théorie de Floquet --- Wikipédia (\url{https://fr.wikipedia.org/wiki/Th%C3%A9or%C3%A8me_de_Floquet})
\item Contexte : Équation de Mathieu --- Wikipédia (\url{https://fr.wikipedia.org/wiki/%C3%89quation_de_Mathieu})
\item Contexte : Oscillateur paramétrique --- Wikipédia (\url{https://fr.wikipedia.org/wiki/Pendule_simple_%C3%A0_r%C3%A9sonance_param%C3%A9trique})
\item Lecture : V. I. Arnold, \textit{Mathematical Methods of Classical Mechanics} (la section sur la résonance paramétrique)
\end{refs}

\bigskip

\begin{problem}[{L'indice d'un champ de vecteurs plan --- Master, short}]
\textbf{ODE-40.} \textbf{Problème.} Soit $ f $ un champ de vecteurs continu sur un ouvert de
$ \mathbb{R}^2 $ et soit $ C $ une courbe fermée simple orientée positivement sur laquelle
$ f $ ne s'annule pas ; on définit l'indice $ I_f(C) $ comme le nombre d'enroulements de
$ f $ le long de $ C $, c'est-à-dire le nombre net de tours complets effectués par le vecteur
$ f $ lorsque le point parcourt $ C $ une fois. Démontrer que $ I_f(C) $ est un entier
inchangé par les déformations continues de $ C $ ne franchissant jamais un zéro de $ f $,
qu'il est égal à la somme des indices des zéros isolés enfermés, qu'un col hyperbolique a pour
indice $ -1 $ tandis qu'un nœud, un foyer ou un centre a pour indice $ +1 $, et enfin que
toute orbite périodique d'un système plan de classe $ C^1 $ a pour indice $ +1 $ --- si bien que
toute orbite périodique doit enfermer des points d'équilibre dont les indices ont pour somme
$ +1 $, et en particulier doit en enfermer au moins un.
\end{problem}

\noindent Poincaré a introduit l'indice dans ses mémoires sur les courbes définies par
des équations différentielles (1881--1886), et c'est l'instrument topologique le moins coûteux du
domaine : une courbe fermée, un entier, et des familles entières de trajectoires sont écartées.
C'est la moitié locale du théorème de Poincaré--Hopf, dont l'énoncé global --- la somme des indices
des zéros d'un champ de vecteurs sur une surface fermée vaut la caractéristique d'Euler --- contient
le théorème de la boule chevelue comme cas de la sphère. Employé avec les arguments de piégeage
d'ODE-14 et le critère négatif d'ODE-28, il est souvent le moyen le plus rapide de décider si un
cycle peut exister.

\begin{refs}
\item Contexte : Théorème de Poincaré-Hopf --- Wikipédia (\url{https://fr.wikipedia.org/wiki/Th%C3%A9or%C3%A8me_de_Poincar%C3%A9-Hopf})
\item Contexte : Indice d'un lacet (nombre d'enroulements) --- Wikipédia (\url{https://fr.wikipedia.org/wiki/Indice_(analyse_complexe)})
\item Lecture : S. H. Strogatz, \textit{Nonlinear Dynamics and Chaos}, chap. 6
\item Lecture : L. Perko, \textit{Differential Equations and Dynamical Systems}, chap. 3
\end{refs}

\bigskip

\begin{problem}[{Couches limites et développements asymptotiques raccordés --- Master}]
\textbf{ODE-41.} \textbf{Problème.} Pour $ \varepsilon > 0 $ petit, on considère le problème aux limites
\[ \varepsilon y'' + (1 + \varepsilon)\,y' + y = 0, \qquad y(0) = 0, \quad y(1) = 1 . \]
\begin{enumerate}
\item[(a)] Le résoudre exactement : factoriser le polynôme caractéristique en
$ (\varepsilon \lambda + 1)(\lambda + 1) $, obtenir
\[ y_\varepsilon(x) = \frac{e^{-x} - e^{-x/\varepsilon}}{e^{-1} - e^{-1/\varepsilon}}, \]
et démontrer que $ y_\varepsilon(x) \to e^{1-x} $ quand $ \varepsilon \to 0^{+} $ pour tout
$ x \in (0,1] $ fixé, alors que $ y_\varepsilon(0) = 0 $ pour tout $ \varepsilon $. Montrer
que la transition se produit sur un intervalle de largeur d'ordre $ \varepsilon $.
\item[(b)] Oublier maintenant la formule et faire ce qu'il faut bien faire pour les équations qui n'en
ont pas. Poser $ \varepsilon = 0 $ donne l'équation \textit{extérieure} du premier ordre
$ y' + y = 0 $, qui ne peut satisfaire les deux conditions au bord ; la résoudre avec la
condition en $ x = 1 $. Changer d'échelle par $ x = \varepsilon X $, établir l'équation
\textit{intérieure} au premier ordre $ Y'' + Y' = 0 $, la résoudre avec $ Y(0) = 0 $, et fixer la
constante restante par raccordement : la limite de la solution intérieure quand
$ X \to \infty $ doit égaler la limite de la solution extérieure quand $ x \to 0^{+} $.
Assembler l'approximation composite $ y \approx e^{1-x} - e^{\,1 - x/\varepsilon} $ comme
extérieure plus intérieure moins partie commune.
\item[(c)] Démontrer que l'approximation est uniformément bonne : en utilisant la solution exacte de
(a), montrer que le supremum sur $ [0,1] $ de la différence entre $ y_\varepsilon $ et la
composite est $ O(\varepsilon) $.
\item[(d)] Expliquer le placement. Démontrer que si l'on postule au contraire une couche en $ x = 1 $
le raccordement est impossible, et identifier le trait de l'équation --- le signe du coefficient de
$ y' $ --- qui décide à quelle extrémité la couche doit se trouver. Puis dire ce que cela a à voir
avec Prandtl : dans un écoulement de viscosité $ \varepsilon $, l'équation extérieure non
visqueuse ne peut satisfaire la condition d'adhérence à la paroi, et la résolution a exactement la
structure calculée ici.
\end{enumerate}

\end{problem}

\noindent Ludwig Prandtl a annoncé la théorie des couches limites lors d'une conférence de
dix minutes à Heidelberg en 1904 : si petite que soit la viscosité, elle ne peut être négligée
dans une mince couche près de la paroi, parce que les équations non visqueuses sont d'ordre
inférieur et ne peuvent imposer la condition d'adhérence. C'est toute la théorie des perturbations
singulières en une phrase --- le petit paramètre multiplie la dérivée d'ordre le plus élevé, si bien
que le problème limite perd une condition au bord, et que la condition perdue réapparaît sous
forme de couche. Friedrichs, Kaplun, Lagerstrom et Van Dyke ont systématisé le raccordement dans
les années 1950 et 1960 ; la même structure gouverne l'approximation WKB près d'un point de
retournement (comparer ODE-29), les variétés lentes de l'oscillation de relaxation (ODE-15), et
l'approximation de l'état quasi stationnaire en cinétique enzymatique.

\begin{refs}
\item Contexte : Développements asymptotiques raccordés --- Wikipédia (\url{https://fr.wikipedia.org/wiki/D%C3%A9veloppements_asymptotiques_raccord%C3%A9s})
\item Contexte : Perturbation singulière --- Wikipédia (\url{https://fr.wikipedia.org/wiki/Perturbation_singuli%C3%A8re})
\item Contexte : Couche limite --- Wikipédia (\url{https://fr.wikipedia.org/wiki/Couche_limite})
\item Lecture : C. M. Bender \& S. A. Orszag, \textit{Advanced Mathematical Methods for Scientists and Engineers}, chap. 9
\end{refs}

\bigskip


\section*{Recherche}

\begin{problem}[{Le seizième problème de Hilbert, seconde partie --- Recherche}]
\textbf{ODE-19.} \textbf{Problème.} Pour un système polynomial plan
$ x' = P(x, y) $, $ y' = Q(x, y) $ de degré $ n $, soit $ H(n) $ le nombre maximal
possible de cycles limites (orbites périodiques isolées). Hilbert a demandé en 1900 la valeur de
$ H(n) $ et les positions possibles des cycles.
(a) Échauffement réalisable dès aujourd'hui : démontrer que $ H(1) = 0 $ --- un système linéaire
plan n'a aucun cycle limite. (Toute orbite périodique d'un système linéaire appartient à un
continuum d'orbites périodiques, donc aucune n'est isolée.)
(b) Problème de compréhension : faire le point sur ce que l'on sait. L'affirmation de Dulac (1923)
selon laquelle chaque système polynomial particulier n'a qu'un nombre fini de cycles limites s'est
révélée lacunaire (Ilyashenko, 1981) et n'a été démontrée qu'ensuite, indépendamment, par Écalle et
Ilyashenko vers 1991--1992. Pourtant aucune borne uniforme n'est connue : même $ H(2) $ est
inconnu --- sa finitude n'est pas démontrée --- alors qu'il existe des systèmes quadratiques à 4 cycles
limites (Shi Songling, 1980). Les minorations de $ H(n) $ croissent comme $ n^2 \ln n $
(Christopher--Lloyd, 1995).
\textit{Statut : ouvert, pour tout $ n \ge 2 $. C'est la seconde partie du seizième problème de
Hilbert et il figure de nouveau sur la liste de Smale pour le XXIe siècle. Une
prépublication de 2024 revendiquant une solution générale a été rapidement réfutée (note de
Buzzi--Novaes, arXiv:2411.09594) --- une leçon vivante sur la nécessité de vérifier un statut plutôt
que de le supposer.}
\end{problem}

\noindent Des 23 problèmes de Hilbert de 1900, celui-ci est parmi les très rares sur
lesquels un siècle de travail n'a produit aucune borne, pas même pour un seul degré. C'est un
problème sur les objets d'ODE-14 --- les cycles limites --- posé uniformément sur tous les champs de
vecteurs polynomiaux, et sa résistance explique pourquoi les résultats de finitude
(Écalle--Ilyashenko) sont célébrés bien qu'ils ne donnent aucun nombre. La synthèse d'Ilyashenko
\og{} Centennial history of Hilbert's 16th problem \fg{} est la carte de référence du champ de bataille.
Voir aussi la page Problèmes ouverts (\url{open-problems.html}) du site.

\begin{refs}
\item Contexte : Seizième problème de Hilbert --- Wikipédia (\url{https://fr.wikipedia.org/wiki/Seizi%C3%A8me_probl%C3%A8me_de_Hilbert})
\item Article : Yu. Ilyashenko, \og{} Centennial history of Hilbert's 16th problem \fg{}, \textit{Bulletin of the AMS} 39 (2002), 301--354
\item Lecture : S. H. Strogatz, \textit{Nonlinear Dynamics and Chaos}, chap. 7 (cycles limites)
\end{refs}

\bigskip

\begin{problem}[{Des singularités sans collision : la conjecture de Painlevé --- Recherche}]
\textbf{ODE-20.} \textbf{Problème.} Une solution du problème newtonien à $ n $ corps présente une
\textit{singularité} à l'instant $ t^* <\infty $ si elle ne peut être prolongée au-delà de
$ t^* $. Les collisions en sont la cause évidente. Painlevé a demandé (1895) : existe-t-il des
singularités \textit{sans} collision --- des corps dont le mouvement devient indéfini en temps fini
alors que deux d'entre eux ne se rencontrent jamais ?
\begin{enumerate}
\item[(a)] Un morceau concret que vous pouvez démontrer : pour deux corps lâchés sans vitesse initiale à
la distance $ d $, établir l'équation de la chute libre radiale et calculer le temps exact
jusqu'à la collision ; en conclure que des singularités de collision se produisent bien. Painlevé
lui-même a démontré que pour $ n = 3 $ \textit{toute} singularité est une collision --- lire
pourquoi l'argument s'arrête à trois corps.
\item[(b)] Problème de compréhension : étudier la résolution. Painlevé conjecturait l'existence de
singularités sans collision pour $ n \ge 4 $. Un tel mouvement est nécessairement non borné :
les corps doivent partir à l'infini en temps fini, en extrayant une énergie non bornée des
quasi-rencontres de binaires qui se resserrent. Mather et McGehee ont construit un mouvement
colinéaire à 4 corps atteignant l'infini en temps fini mais utilisant des collisions binaires
(1975) ; Zhihong (Jeff) Xia a construit une véritable singularité sans collision pour le problème
spatial à 5 corps (annoncée dans sa thèse de 1988 ; \textit{Annals of Mathematics}, 1992) ; Jinxin
Xue a réglé le dernier cas, le problème plan à 4 corps (annoncé en 2014 ; \textit{Acta Mathematica},
2020). Expliquer le mécanisme de Xia --- le corps oscillant faisant la navette entre deux binaires ---
et pourquoi $ n = 3 $ l'interdit.
\textit{Statut : la conjecture de Painlevé est résolue --- des singularités sans collision existent pour
tout $ n \ge 4 $ (Xia 1992 pour $ n \ge 5 $ ; Xue 2020 pour $ n = 4 $). Les questions
quantitatives (mesure de ces données initiales, que l'on croit nulle ; classification des
singularités) restent ouvertes.}
\end{enumerate}

\end{problem}

\noindent Painlevé a posé la question dans ses leçons de Stockholm en 1895, et il a fallu
presque un siècle pour y répondre. Les constructions sont des chefs-d'œuvre de la théorie
qualitative vers laquelle cette page tend : aucune formule, mais un contrôle sans faille des
orbites au travers d'une infinité de quasi-collisions. L'article de Xue de 2020 --- 135 pages dans
\textit{Acta Mathematica} --- qui clôt le cas $ n = 4 $ rappelle que la mécanique classique livre
encore de durs théorèmes modernes.

\begin{refs}
\item Contexte : Conjecture de Painlevé --- Wikipédia (\url{https://fr.wikipedia.org/wiki/Conjecture_de_Painlev%C3%A9})
\item Article : Z. Xia, \og{} The existence of noncollision singularities in Newtonian systems \fg{}, \textit{Annals of Mathematics} 135 (1992), 411--468
\item Article : J. Xue, \og{} Non-collision singularities in a planar 4-body problem \fg{}, \textit{Acta Mathematica} 224 (2020), 253--388, arXiv:1409.0048 (\url{https://arxiv.org/abs/1409.0048})
\end{refs}

\bigskip

\begin{problem}[{Le Système solaire est-il stable ? --- Recherche}]
\textbf{ODE-21.} \textbf{Problème.} La plus ancienne question de la dynamique : les planètes restent-elles à peu
près sur leurs orbites actuelles pour toujours, ou l'accumulation lente des perturbations
gravitationnelles peut-elle éjecter une planète ou provoquer une collision ?
\begin{enumerate}
\item[(a)] Un noyau soluble : démontrer que le problème à deux corps (Kepler) est intégrable --- établir
la conservation de l'énergie et du moment cinétique, et montrer que le vecteur de
Laplace--Runge--Lenz $ A = p \times L - \frac{G m_1 m_2\, m\, \hat{r}}{1} $ (convenablement
normalisé) est une constante du mouvement ; en déduire que les orbites bornées sont des ellipses
sans précession.
\item[(b)] Problème de compréhension : la théorie KAM. Énoncer informellement le théorème de
Kolmogorov--Arnold--Moser (1954--1963) --- sous une petite perturbation d'un système intégrable, la
plupart des tores quasi périodiques survivent --- et expliquer à la fois ce que dit l'application
d'Arnold à propos d'un système planétaire aux masses planétaires suffisamment minuscules, et
pourquoi les masses réelles de notre Système solaire se situent bien au-delà des hypothèses
rigoureuses de petitesse.
\item[(c)] Problème de compréhension : le verdict numérique. Les intégrations à long terme de Laskar
(à partir de 1989) montrent que le Système solaire interne est chaotique, avec un temps de
Liapounov d'environ 5 millions d'années, et Laskar \& Gastineau (2009) ont trouvé une
probabilité d'environ 1 \% que l'orbite de Mercure se déstabilise pendant la durée de vie restante
du Soleil. Expliquer ce que signifie un temps de Liapounov pour la question, et ce que de telles
simulations peuvent et ne peuvent pas démontrer.
\textit{Statut : ouvert en tant que problème rigoureux. La stabilité de type KAM n'est démontrée que
pour des masses irréalistement petites ; les données numériques indiquent un chaos marginal avec
une probabilité faible mais non nulle de catastrophe. Un théorème couvrant le Système solaire réel
n'existe pas.}
\end{enumerate}

\end{problem}

\noindent Newton pensait que Dieu pourrait devoir intervenir de temps à autre pour
remettre les planètes en place ; Laplace et Lagrange ont démontré la stabilité au premier ordre en
les masses, et le mémoire de Poincaré d'ODE-18 a brisé l'espoir que les ordres supérieurs se
tiendraient bien. La théorie KAM --- l'une des plus profondes réussites de l'analyse du
XXe siècle --- a rétabli la stabilité \og{} pour la plupart des conditions initiales \fg{} des
systèmes quasi intégrables, mais entre ses hypothèses et le Système solaire réel s'ouvre un
fossé que personne n'a comblé. La question qui a fondé les systèmes dynamiques est toujours ouverte
en son centre.

\begin{refs}
\item Contexte : Stabilité du Système solaire --- Wikipédia (en anglais) (\url{https://en.wikipedia.org/wiki/Stability_of_the_Solar_System})
\item Contexte : Théorème KAM --- Wikipédia (\url{https://fr.wikipedia.org/wiki/Th%C3%A9or%C3%A8me_KAM})
\item Article : J. Laskar \& M. Gastineau, \og{} Existence of collisional trajectories of Mercury, Mars and Venus with the Earth \fg{}, \textit{Nature} 459 (2009), 817--819
\item Lecture : J. Moser, \textit{Stable and Random Motions in Dynamical Systems} (1973)
\end{refs}

\bigskip

\begin{problem}[{La conjecture de Markus--Yamabe --- Recherche, short}]
\textbf{ODE-31.} \textbf{Problème.} Markus et Yamabe ont conjecturé en 1960 que si
$ f : \mathbb{R}^n \to \mathbb{R}^n $ est de classe $ C^1 $ avec $ f(0) = 0 $ et que toute
valeur propre de la jacobienne $ Df(x) $ est de partie réelle négative en \textit{tout} $ x $,
alors l'origine est globalement asymptotiquement stable pour $ x' = f(x) $. Démontrer vous-même
le cas linéaire, puis rendre compte avec soin de la résolution --- la conjecture est un théorème
dans le plan (démontré indépendamment par Fessler, Gutiérrez et Glutsyuk vers 1993--1995) et
fausse en toute dimension $ n \ge 3 $, où Cima, van den Essen, Gasull, Hubbers et Mañosas (1997)
ont exhibé un contre-exemple polynomial à orbite fuyante.
\end{problem}

\noindent La conjecture est l'espoir naturel qu'une hypothèse vérifiée en tout point
doive coudre la stabilité locale en une stabilité globale, et cet espoir est exactement juste en
dimension deux et exactement faux au-dessus. Le contre-exemple est issu de l'école de Nimègue
travaillant sur la conjecture jacobienne, et son champ de vecteurs --- aujourd'hui appelé système de
Cima --- est assez explicite pour être vérifié à la main. \textit{Statut : réglé ; vrai pour $ n = 2 $,
faux pour $ n \ge 3 $. Les variantes (non autonome, en dimension infinie, discrète) restent
actives.}

\begin{refs}
\item Contexte : Conjecture de Markus--Yamabe --- Wikipédia (en anglais) (\url{https://en.wikipedia.org/wiki/Markus%E2%80%93Yamabe_conjecture})
\item Article : A. Cima, A. van den Essen, A. Gasull, E. Hubbers \& F. Mañosas, \og{} A polynomial counterexample to the Markus--Yamabe conjecture \fg{}, \textit{Advances in Mathematics} 131 (1997), 453--457
\item Lecture : L. Perko, \textit{Differential Equations and Dynamical Systems}, chap. 2
\end{refs}

\bigskip

\begin{problem}[{Le sixième problème de Smale : compter les configurations centrales --- Recherche, short}]
\textbf{ODE-32.} \textbf{Problème.} Un \textit{équilibre relatif} du problème newtonien à
$ n $ corps est une configuration plane qui tourne rigidement autour de son centre de masse ---
de façon équivalente, une configuration où l'accélération de chaque corps est le même multiple
négatif de son vecteur position. Le sixième problème de Smale pour le XXIe siècle
demande si, pour tout $ n $ et tout choix de masses positives, le nombre de telles
configurations est fini à rotation et homothétie près ; démontrer le cas $ n = 3 $, où les trois
configurations colinéaires d'Euler (1767) et celle, équilatérale, de Lagrange (1772) sont les
seules, puis rendre compte du statut moderne.
\end{problem}

\noindent La question est un problème déguisé de géométrie algébrique réelle : les
configurations sont les solutions réelles d'un système polynomial, et la finitude n'échoue que si
toute une courbe de solutions apparaît. Hampton et Moeckel ont démontré la finitude pour
$ n = 4 $ en 2006 par calcul symbolique exact, encadrant le compte entre 32 et 8472 ; Albouy et
Kaloshin ont réglé $ n = 5 $ en 2012, sauf pour des masses appartenant à une sous-variété de
codimension deux. \textit{Statut : ouvert pour tout $ n \ge 6 $ (et pour les masses exceptionnelles à
cinq corps). Voir la page Problèmes ouverts (\url{open-problems.html}) du site pour la liste
de Smale.}

\begin{refs}
\item Contexte : Configuration centrale --- Wikipédia (en anglais) (\url{https://en.wikipedia.org/wiki/Central_configuration})
\item Contexte : Problèmes de Smale --- Wikipédia (\url{https://fr.wikipedia.org/wiki/Probl%C3%A8mes_de_Smale})
\item Article : M. Hampton \& R. Moeckel, \og{} Finiteness of relative equilibria of the four-body problem \fg{}, \textit{Inventiones Mathematicae} 163 (2006), 289--312
\item Article : A. Albouy \& V. Kaloshin, \og{} Finiteness of central configurations of five bodies in the plane \fg{}, \textit{Annals of Mathematics} 176 (2012), 535--588
\end{refs}

\bigskip

\begin{problem}[{Diffusion d'Arnold : les actions peuvent-elles dériver ? --- Recherche}]
\textbf{ODE-33.} \textbf{Problème.} On écrit un hamiltonien quasi intégrable en variables action--angle
\[ H(I, \theta) = H_0(I) + \varepsilon H_1(I, \theta), \qquad I \in \mathbb{R}^n, \ \theta \in \mathbb{T}^n . \]
La théorie KAM (ODE-21) affirme que la plupart des tores invariants survivent pour
$ \varepsilon $ petit. Arnold a demandé en 1964 si les actions peuvent néanmoins errer sur une
distance d'ordre un.
\begin{enumerate}
\item[(a)] Point d'entrée : démontrer que pour $ \varepsilon = 0 $ les actions sont exactement
constantes, de sorte que toute dérive est un effet de perturbation. Démontrer ensuite que pour
$ n = 2 $ une dérive d'ordre un est \textit{impossible} à petit $ \varepsilon $ : les tores KAM
survivants sont de dimension deux à l'intérieur d'un niveau d'énergie de dimension trois et le
séparent donc, piégeant toute trajectoire entre deux tores voisins.
\item[(b)] Étudier l'exemple d'Arnold de 1964, le premier système où les actions bougent de façon
démontrée. Décrire sa chaîne de tores partiellement hyperboliques (\og{} à moustaches \fg{}) dont les
variétés stables et instables se coupent transversalement, et expliquer comment un argument de
pistage produit une orbite unique suivant toute la chaîne, changeant son action de
$ O(1) $.
\item[(c)] Expliquer pourquoi la dérive doit être lente : énoncer le théorème de Nekhorochev (1977), qui
borne la variation des actions par une petite puissance de $ \varepsilon $ sur des temps
exponentiellement longs en une puissance de $ 1/\varepsilon $ pour $ H_0 $ escarpé, et le
concilier avec (b).
\item[(d)] Problème de compréhension : rendre compte honnêtement de l'état de l'art. La diffusion est
établie dans les cadres \textit{a priori instables} (Chierchia--Gallavotti ;
Delshams--de la Llave--Seara ; Treschev ; Cheng--Yan), et la monographie de Kaloshin et Zhang de 2020
la démontre pour des perturbations typiques à deux degrés et demi de liberté, achevant un
programme annoncé par John Mather en 2003. Dire précisément ce que \og{} typique \fg{} signifie dans ces
théorèmes et ce qu'exigerait une solution complète.
\textit{Statut : ouvert en général. La conjecture selon laquelle la diffusion d'Arnold est générique
dans les systèmes quasi intégrables à $ n \ge 3 $ degrés de liberté n'est pas démontrée ; les
résultats généraux annoncés sont encore en cours d'assimilation par la communauté.}
\end{enumerate}

\end{problem}

\noindent La note de trois pages d'Arnold en 1964 est l'un des courts articles les plus
lourds de conséquences de la dynamique : elle a montré que les îlots de stabilité de KAM, si
abondants soient-ils, n'enferment pas nécessairement une trajectoire dès qu'il y a plus de deux
degrés de liberté, parce que les interstices entre tores communiquent. Le phénomène est
physiquement réel et physiquement lent --- on l'invoque pour le comportement à long terme des
orbites d'astéroïdes et pour la perte de confinement des faisceaux dans les accélérateurs de
particules --- et l'estimation de Nekhorochev explique pourquoi le Système solaire d'ODE-21 peut
être chaotique sans s'être disloqué. Le sujet est l'un des rares où un mécanisme a été décrit
avant que quiconque puisse démontrer qu'il se produit génériquement, et le démontrer occupe les
chercheurs depuis soixante ans.

\begin{refs}
\item Contexte : Diffusion d'Arnold --- Wikipédia (en anglais) (\url{https://en.wikipedia.org/wiki/Arnold_diffusion})
\item Contexte : Théorème KAM --- Wikipédia (\url{https://fr.wikipedia.org/wiki/Th%C3%A9or%C3%A8me_KAM})
\item Article : V. I. Arnold, \og{} Instability of dynamical systems with several degrees of freedom \fg{}, \textit{Soviet Mathematics Doklady} 5 (1964), 581--585
\item Lecture : V. Kaloshin \& K. Zhang, \textit{Arnold Diffusion for Smooth Systems of Two and a Half Degrees of Freedom}, Annals of Mathematics Studies 208 (Princeton, 2020)
\end{refs}

\bigskip

\begin{problem}[{La conjecture de Wright et la mémoire d'une équation à retard --- Recherche}]
\textbf{ODE-42.} \textbf{Problème.} L'équation de Wright est l'équation différentielle scalaire à retard
\[ y'(t) = -\alpha\, y(t-1)\,\bigl[ 1 + y(t) \bigr], \qquad \alpha > 0, \]
dont la donnée initiale est une fonction continue sur $ [-1, 0] $ avec $ 1 + y > 0 $.
\begin{enumerate}
\item[(a)] Montrer d'où elle vient : en partant de l'équation logistique retardée de Hutchinson
$ N'(t) = r N(t)\bigl[ 1 - N(t-T)/K \bigr] $ pour une population qui ne réagit à
l'encombrement qu'après un délai $ T $, substituer $ y = N/K - 1 $ et changer l'échelle de
temps pour obtenir l'équation de Wright avec $ \alpha = rT $. Expliquer ce que signifie
biologiquement la contrainte $ 1 + y > 0 $.
\item[(b)] Linéariser en $ y \equiv 0 $ pour obtenir $ y' = -\alpha\,y(t-1) $, substituer
$ y = e^{\lambda t} $, et étudier l'équation caractéristique $ \lambda = -\alpha e^{-\lambda} $.
Démontrer qu'elle a une infinité de racines complexes --- l'équation linéaire a donc une infinité de
solutions exponentielles indépendantes, et aucun système fini d'équations différentielles
ordinaires ne peut la reproduire --- et démontrer que toute racine est de partie réelle négative
exactement lorsque $ \alpha < \pi/2 $, une paire conjuguée traversant l'axe imaginaire en
$ \lambda = \pm i\pi/2 $ lorsque $ \alpha = \pi/2 $.
\item[(c)] Énoncer la conjecture de Wright : pour tout $ \alpha \in (0, \pi/2] $, l'origine attire
toute solution. Expliquer soigneusement pourquoi (b) ne la démontre pas --- la stabilité
asymptotique locale ne dit rien des données éloignées de zéro --- et lire l'argument de Wright
lui-même de 1955, qui règle $ \alpha \le 3/2 $, pour voir ce qu'une démonstration globale doit
contrôler.
\item[(d)] Rendre compte de la résolution. Jan Bouwe van den Berg et Jonathan Jaquette ont refermé la
fenêtre restante jusqu'à $ \pi/2 $ en 2018 par une analyse minutieuse du voisinage de la
bifurcation de Hopf en $ \alpha = \pi/2 $, complétant des travaux antérieurs qui couvraient les
valeurs de paramètre plus éloignées de la bifurcation ; le même cercle d'idées, dans l'article de
Jaquette de 2019, a démontré la conjecture de Jones de 1962, selon laquelle pour
$ \alpha > \pi/2 $ la solution périodique à oscillation lente est unique à translation
temporelle près. Expliquer comment la plage de paramètres se partage entre l'analyse et les
encadrements rigoureux assistés par ordinateur, et dire précisément quelle étape devrait être
remplacée pour obtenir une démonstration à la main.
\textit{Statut : les deux conjectures sont des théorèmes --- celle de Wright (van den Berg--Jaquette,
2018) et celle de Jones (Jaquette, 2019, achevant des résultats partiels antérieurs). Les deux
démonstrations sont assistées par ordinateur ; une démonstration purement analytique de l'une ou
l'autre est un problème ouvert, et pour les équations scalaires à retard en général, décider de la
stabilité globale reste hors de portée.}
\end{enumerate}

\end{problem}

\noindent Les équations à retard sont ce que l'on obtient chaque fois que la réponse
suit la cause avec un délai : les modèles proie--prédateur héréditaires de Volterra en 1928,
l'analyse par Minorsky des stabilisateurs de navires dans les années 1940, la logistique retardée
de Hutchinson en 1948, le modèle de cellules sanguines de Mackey et Glass (1977) qui a fait d'une
équation scalaire à retard un exemple standard de chaos, et toute boucle d'asservissement à
capteur retardé. L'état d'une telle équation est une fonction entière sur un intervalle, si bien
que l'espace des phases est de dimension infinie et que le flot n'est pas inversible --- ce qui
explique pourquoi une équation scalaire à retard peut faire ce qu'aucune équation différentielle
ordinaire scalaire ne peut, et pourquoi l'équation d'apparence innocente de Wright a résisté
soixante ans. Qu'elle soit finalement tombée sous un mélange d'analyse des bifurcations et de
calcul numérique validé la place aux côtés du théorème de Tucker sur l'attracteur de Lorenz
(ODE-17) comme cas d'étude de ce qu'est devenue la démonstration assistée par ordinateur.

\begin{refs}
\item Contexte : Équation différentielle à retard --- Wikipédia (\url{https://fr.wikipedia.org/wiki/%C3%89quation_diff%C3%A9rentielle_%C3%A0_retard})
\item Article : E. M. Wright, \og{} A non-linear difference-differential equation \fg{}, \textit{Journal für die reine und angewandte Mathematik} 194 (1955), 66--87
\item Article : J. B. van den Berg \& J. Jaquette, \og{} A proof of Wright's conjecture \fg{}, \textit{Journal of Differential Equations} 264 (2018), 7412--7462, arXiv:1704.00029 (\url{https://arxiv.org/abs/1704.00029})
\item Article : J. Jaquette, \og{} A proof of Jones' conjecture \fg{}, \textit{Journal of Differential Equations} 266 (2019), 3818--3859, arXiv:1801.09806 (\url{https://arxiv.org/abs/1801.09806})
\end{refs}

\bigskip

\begin{problem}[{La conjecture de stabilité --- Recherche, short}]
\textbf{ODE-43.} \textbf{Problème.} Un champ de vecteurs de classe $ C^1 $ sur une variété
fermée est \textit{structurellement stable} si tout champ suffisamment $ C^1 $-proche lui est
topologiquement équivalent --- même portrait de phase à un homéomorphisme envoyant orbites sur
orbites près --- et Palis et Smale ont conjecturé vers 1970 que cela se produit exactement lorsque le
champ satisfait l'axiome A et la condition de transversalité forte. Démontrer l'échauffement plan
avec le critère d'Andronov--Pontriaguine --- montrer qu'un champ sur le disque possédant un point
d'équilibre non hyperbolique, ou une connexion de col à col, admet des perturbations
arbitrairement petites modifiant son portrait de phase --- puis rendre compte avec soin du cas
général, en disant exactement ce qui a été démontré et dans quelle topologie.
\end{problem}

\noindent Andronov et Pontriaguine ont introduit les \og{}~systèmes
grossiers~\fg{} en 1937 et les ont caractérisés dans le plan ; Peixoto a démontré vers 1962
que sur une surface compacte orientable les champs structurellement stables sont ouverts et
denses, si bien qu'en dimension deux le système typique est un système que l'on peut dessiner.
Smale espérait qu'il en irait de même en dimension supérieure et l'a lui-même réfuté dans les
années 1960 --- les enchevêtrements homoclines sont indéracinables --- après quoi la stabilité
structurelle est devenue une propriété à caractériser plutôt qu'à espérer. La suffisance de
l'axiome A et de la transversalité forte a été démontrée par Robbin (1971) et Robinson (1976) ; la
réciproque en topologie $ C^1 $ fut le théorème célèbre de Ricardo Mañé pour les
difféomorphismes (1987), étendu aux flots par Shuhei Hayashi (1997) grâce à son lemme de
connexion. \textit{Statut : réglé en topologie $ C^1 $ ; ouvert pour $ C^r $ avec $ r \ge 2 $, où
les lemmes de perturbation qui portent toutes les démonstrations connues n'ont aucun
substitut.}

\begin{refs}
\item Contexte : Stabilité structurelle --- Wikipédia (en anglais) (\url{https://en.wikipedia.org/wiki/Structural_stability})
\item Contexte : Critère d'Andronov--Pontriaguine --- Wikipédia (en anglais) (\url{https://en.wikipedia.org/wiki/Andronov%E2%80%93Pontryagin_criterion})
\item Article : R. Mañé, \og{} A proof of the $ C^1 $ stability conjecture \fg{}, \textit{Publications Mathématiques de l'IHÉS} 66 (1987), 161--210
\item Article : S. Hayashi, \og{} Connecting invariant manifolds and the solution of the $ C^1 $ stability and $\Omega$-stability conjectures for flows \fg{}, \textit{Annals of Mathematics} 145 (1997), 81--137 (rectificatif, \textit{Annals of Mathematics} 150 (1999), 353--356)
\end{refs}

\bigskip


\section*{Pour aller plus loin}


\vfill
\noindent\emph{Hors de la Caverne} --- des probl\`emes, pas de r\'eponses.
\end{document}
